% !TEX TS-program = pdflatex
% !BIB TS-program = biber
% This file should be compiled with pdfLaTeX + biber.
% Run: pdflatex main.tex && biber main && pdflatex main.tex && pdflatex main.tex


\documentclass[11pt]{article}

% UTF-8 encoding and classic babel (for pdflatex)
\usepackage[utf8]{inputenc}
\usepackage[T1]{fontenc}
\usepackage{textcomp}
\usepackage{upquote}
\usepackage[a4paper,top=2cm,bottom=2cm,left=2cm,right=2cm,marginparwidth=1.75cm]{geometry}

% Arial via Helvetica clone (pdflatex-compatible), per IU formatting guidelines
\usepackage[scaled]{helvet}
\renewcommand{\familydefault}{\sfdefault}

% 1.5 line spacing
\usepackage{setspace}
\onehalfspacing

% No first-line indent, 6pt spacing after paragraphs
\usepackage{parskip}
\setlength{\parskip}{6pt}
\setlength{\parindent}{0pt}

% Heading sizes/spacing per IU guidelines: L1 16pt (12pt before/after), L2 14pt (12pt before, 6pt after), L3 11pt (12pt before, 6pt after)
\usepackage{titlesec}
\titleformat{\section}{\normalfont\fontsize{16}{19}\bfseries}{\thesection}{1em}{}
\titlespacing*{\section}{0pt}{12pt}{12pt}
\titleformat{\subsection}{\normalfont\fontsize{14}{17}\bfseries}{\thesubsection}{1em}{}
\titlespacing*{\subsection}{0pt}{12pt}{6pt}
\titleformat{\subsubsection}{\normalfont\fontsize{11}{13}\bfseries}{\thesubsubsection}{1em}{}
\titlespacing*{\subsubsection}{0pt}{12pt}{6pt}

\usepackage[backend=biber, style=apa]{biblatex}
\addbibresource{references.bib}

% Math packages
\usepackage{intcalc}
\usepackage{amsmath, amssymb, amsthm, amsfonts, mathtools, bm}
\usepackage{siunitx}
\usepackage{physics}
\usepackage[version=4]{mhchem}
\usepackage{chemformula}
\usepackage{tensor}
\usepackage{hepnames}
\usepackage{braket}

% Tables and figures
\usepackage{adjustbox}
\usepackage{tabularx}
\usepackage{graphicx}
\usepackage{float}
\usepackage[ruled,vlined]{algorithm2e}

% Misc
\usepackage{caption}
\usepackage{csquotes}

% Language support
\usepackage[english]{babel}

% Layout and links
\usepackage[colorlinks=true, allcolors=blue]{hyperref}


\begin{document}

% ===== Title page (page I, no number shown, not listed in TOC) =====
\pagenumbering{roman}
\setcounter{page}{1}
\thispagestyle{empty}
\begin{titlepage}
    \centering
    \vspace*{2cm}
    {\fontsize{16}{19}\bfseries Algorithmic Addiction in Social Recommender Systems: Development of an Ethical AI Audit and Pro-Social Pivot Strategy under the EU Digital Fairness Act\par}
    \vspace{1.5cm}
    {\large CAPSTONE PROJECT (CONSULTING REPORT)\par}
    \vspace{1cm}
    \begin{tabular}{ll}
        Institution: & IU International University of Applied Sciences \\[4pt]
        Degree Programme: & Master of Business Administration (MBA), 90 ECTS, LSBU Dual-Award \\[4pt]
        Author: & Mohit Nirwan \\[4pt]
        Matriculation Number: & IU14087991 \\[4pt]
        Supervisor/Tutor: & Dr Tobias Broweleit \\[4pt]
        Submission Date: & [INSERT SUBMISSION DATE] \\
    \end{tabular}
\end{titlepage}

\setcounter{page}{2}

% ===== Executive Summary (front matter, per Capstone Guidelines: ~200-250 words + keywords) =====
\section*{Executive Summary}
\addcontentsline{toc}{section}{Executive Summary}
This executive summary concisely frames the client problem, methodological approach, principal findings, and actionable recommendations. Problem: major social platforms exhibit algorithmically mediated engagement dynamics that create elevated addiction risks for users and expose firms to intensifying regulatory threat, particularly from EU/DFA policy developments. Approach: we conducted a secondary qualitative synthesis of existing studies, constructed a comparative coding matrix to surface cross-platform mechanisms, and developed an Audit Toolkit that turns assessment of algorithmic harms into a repeatable process for product, policy and legal teams. Key findings: (1) a systemic engagement-addiction risk is observable across TikTok, Instagram and LinkedIn-type platforms, driven by recommendation loops, short-form reinforcement and reward-timing; (2) these mechanisms are technologically diffuse but functionally similar, making platform-specific fixes insufficient; and (3) regulatory momentum is clear: the EU and the anticipated Digital Fairness Act signal a direction of travel toward stricter oversight of personalization and monetization practices. Headline recommendations: adopt the Ethical AI Audit Toolkit immediately to identify high-risk features and remediation pathways; execute a Pro-Social AI pivot that prioritizes attention stewardship and positive social outcomes in model objectives; and realign monetization with wellbeing by decoupling ad incentives from engagement-maximization, introducing alternative revenue models and embedding user wellbeing KPIs into executive compensation. These steps mitigate legal exposure and preserve long-term user trust and platform value.

\section*{Keywords}
\addcontentsline{toc}{section}{Keywords}
Keywords: algorithmic addiction, social recommender systems, dark patterns, Digital Fairness Act, ethical AI audit, pro-social design, EU platform regulation.

\newpage
\tableofcontents
\newpage

% List of Figures omitted: this report contains 0 figures, below the 3-figure threshold in the IU Formatting Guidelines.
\listoftables
\addcontentsline{toc}{section}{List of Tables}

\section*{List of Abbreviations}
\addcontentsline{toc}{section}{List of Abbreviations}
\begin{itemize}
\item AI -- Artificial Intelligence
\item AIAS-5 -- Algorithmic Impact Assessment Scale
\item CI/CD -- Continuous Integration / Continuous Deployment
\item CSR -- Corporate Social Responsibility
\item DFA -- Digital Fairness Act
\item DSA -- Digital Services Act
\item HITL -- Human-in-the-Loop
\item SDT -- Self-Determination Theory
\item UX -- User Experience
\item VLOP -- Very Large Online Platform
\end{itemize}

\newpage
\pagenumbering{arabic}
\setcounter{page}{1}


\section{Problem Description and Regulatory Environment}

\subsection{Motivation and Practical Relevance}
\label{sec:motivation_practical_relevance} Dark-pattern research has documented the anatomy of manipulative design in striking detail. Subscription flows are engineered to make cancelling harder than signing up. Illusory-urgency countdowns and misleading brand associations are built into standard UI templates \parencites[390]{7896734}. This is deliberate design: behavioural evidence confirms that a "only two left" message shifts a neutral non-purchase into a perceived loss \parencites[135]{8745065}, because scarcity cues activate loss-framing over deliberation. As AI-driven personalisation scales these techniques, the exploitation of predictable cognitive vulnerabilities becomes systematic, making a targeted audit methodology practically necessary.

European regulation is already moving to curb this trajectory. The DSA requires very large platforms to assess "systemic risks" to civic discourse and electoral processes, backed by transparency reports and non-profiled recommendation options \parencites[12]{6502223}[2]{9360054}. The AI Act goes further, banning subliminal manipulation and manipulative effects that significantly harm behaviour and placing such practices in the "unacceptable risk" tier \parencites[12]{6502223}[390]{7896734}. A planned Digital Fairness Act would regulate addictive design and unfair personalisation as a distinct problem space bridging the DSA and consumer protection law. The regulatory direction is clear even before that Act is codified.

Global AI-governance proposals reinforce the same logic. Kostenko's AI Law Model requires systems to respect human dignity, autonomy and psycho-emotional comfort \parencites[61--62]{2977039}, prohibiting designs that manipulate behaviour or cause digital addiction. The fairness principle bans systems built on distorted data or that consolidate social inequality, with certification, suspension powers and public ethics-impact reports mandated for high-risk domains \parencites[73--74]{2977039}.

Digital-literacy gaps sharpen the urgency. Survey evidence on GDPR/DSA/AI Act awareness shows high self-assessed basic skills but persistent difficulty recognising hidden content and sophisticated manipulation \parencites[2]{9360054}[12]{6502223}. Users are sceptical that regulation alone guarantees safety, and individual awareness cannot carry the protective burden alone. OECD and UNESCO AI-ethics instruments articulate fairness and transparency norms \parencites[9]{7461574}[12]{6502223}[12]{8367105}, yet the "principles-to-practice gap" literature finds high-level declarations rarely specify operational procedures for data governance or post-deployment monitoring \parencites[9]{7461574}[12]{8367105}. Accountability registries and human-in-the-loop oversight are becoming deployment prerequisites for high-risk sectors \parencites[83]{2977039}[11]{5419804}[21]{6805789}[108]{2977039}. Social recommender systems shaping attention at scale belong in the same category. An audit and pro-social pivot strategy that is both ethics-grounded and transparent is, on those grounds, academically necessary and regulatorily urgent.

\subsection{Research Objectives and Central Questions}
\label{sec:research_objectives_central_questions} Interface tactics in digital retail are not isolated creative decisions. Research on urgency marketing and dark patterns shows they are systematically organised around psychological vulnerabilities and short-term revenue objectives \parencites[136]{8745065}. This project's first objective is a secondary-research synthesis mapping such manipulation-oriented choices in recommender-driven environments through a coding matrix drawing on scarcity, psychological targeting and revenue optimisation constructs.

A second objective follows regulatory ambition. The AI Act bans subliminal manipulation and manipulative effects significantly harming behaviour; the DSA requires very large platforms to assess systemic risks to civic discourse and to offer at least one non-profiled recommendation option \parencites[2248]{6502223}. A forthcoming Digital Fairness Act will complement this by targeting addictive design and unfair personalisation as its own problem space \parencites[390]{7896734}. Translating these obligations into operational audit questions is the practical task. Here the DFA is policy trajectory; the AI Act and DSA remain the binding anchor frameworks \parencites[390]{7896734}[2248]{6502223}.

A third objective is methodological. Qualitative and secondary designs examining AI systems and regulatory texts stress triangulating regulatory intent, platform practice and documented harm \parencites[7]{6295677}[6]{6502223}[3]{5142478}. This study works from interpretive synthesis of secondary materials; proprietary telemetry was not available. Black-box algorithms are assessed through documented UI behaviour and regulator findings, following Garr's convergent-parallel content-curation design and Chakraborty et al.'s mixed-methods education-ethics work as templates for evaluating algorithmic risk, calibrated to that same opacity, without internal model access \parencites[3]{5142478}[5]{1308731}.

The central questions operate at the level of observable outputs and regulatory alignment. Do engagement-optimising recommendation architectures embed structural incentives similar to those identified in political recommender systems, where optimisation objectives and training dynamics contribute to polarisation \parencites[6]{6502223}? How far do existing and emerging EU instruments (AI Act, DSA, Digital Markets Act, projected DFA) function as a de facto governance regime for addictive design \parencites[390]{7896734}[2248]{6502223}[2]{1456534}?

A final objective draws on ethical AI marketing and corporate digital responsibility research. Asante and Adeoba \parencites[11]{8200246}[16]{8200246} identify Communities of Practice and CSR as mechanisms turning internal ethical learning into policies and algorithm audits, with stakeholder trust as the core outcome. Catapang's Ethics-by-Design account grounds human agency, fairness and transparency in specific lifecycle tasks, moving beyond abstract principles \parencites[12]{8367105}. The working objective here is a practical audit toolkit and decision matrix, usable by management teams, linked to CSR governance and oriented toward sustaining engagement within EU risk-based regulation \parencites[11]{8200246}[12]{8367105}[4]{1456534}. Closing that principles-to-practice gap is the benchmark against which this study's recommendations are judged.

\subsection{Scope, Boundaries and Key Assumptions}
\label{sec:scope_boundaries_assumptions} This project draws on qualitative secondary research: literature-based inquiry and reflexive thematic analysis synthesise academic, policy and industry documents \parencites[8]{8200246}[6]{6502223}, treating publications and regulatory texts as primary data and building understanding inductively, without primary empirical measurement or direct access to proprietary systems. Claims about algorithmic addiction and manipulative interface patterns are mediated through documented analyses of recommendation architectures and regulatory regimes; no fresh behavioural experiments were conducted.

Black-box AI systems are examined through observable outputs and described objective functions; trade secrecy rules out source-code inspection and limits detailed technical analysis. This thesis adopts the same boundary as scholars studying democratic harms from recommendation architectures, framing audits around externally assessable properties: explainability obligations, independent audit requirements and risk classifications \parencites[6]{6502223}[16]{6502223}[244]{2977039}. Internal weights are never reconstructed; every analytical conclusion about algorithmic risk derives from published regulatory findings and academic syntheses instead.

The regulatory frame draws on formally adopted EU instruments (DSA, AI Act) alongside extended AI law proposals specifying transparency, non-discrimination and audit obligations for high-risk systems, including algorithmic passports and national AI registers \parencites[244]{2977039}[243]{2977039}. The AI Law Model's digital-neutrality principle prohibits information-flow distortion and hidden censorship, and requires documented decision logic and independent audit with violations recorded in a public Register of Digital Neutrality Incidents \parencites[109]{2977039}. Any forward-looking reference to a Digital Fairness Act \parencites[391]{7896734} is treated as policy direction anchored in existing materials; it is not yet codified law binding on current platform practice.

Organisational governance scope emphasises Communities of Practice, CSR mechanisms and AI ethics committees translating shared ethical learning into policies and accountability structures \parencites[11]{8200246}[16]{8200246}. Ethical AI marketing research \parencites[11]{8200246} restricts the governance focus to transparency, privacy, fairness and user control as the primary trust drivers, leaving technical compliance metrics outside scope.

Three key assumptions underpin the analysis. First, digital-literacy and regulatory-awareness surveys \parencites[2]{9360054} indicate persistent user vulnerability despite high self-rated competence. Second, human-in-the-loop oversight remains indispensable: regulatory texts call for supervision channels and appeal procedures, and educational HITL models show combining machine efficiency with human judgment sustains trust \parencites[244]{2977039}[21]{6805789}[3]{5683952}. Third, independent auditing and transparency are enforceable governance levers: AI law proposals establish mandatory audits and registration, and democratic recommendation governance work advocates annual third-party audits against content-policy and diversity benchmarks \parencites[16]{6502223}[244]{2977039}. The Toolkit this study develops therefore operates within published regulatory expectations and observable system behaviours; it makes no claims that require access to proprietary model weights or internal platform data.

\subsection{Conceptual Clarification of Algorithmic Addiction}
\label{sec:conceptual_algorithmic_addiction} Short-video platforms offer the clearest empirical entry point. Neural-network recommendation engines continuously optimise content matching, using micro-duration stimuli and variable reward schedules to trigger dopamine release and induce an "auto-pilot mode" of persistent scrolling \parencites[2]{8958037}, already scrutinised under the DSA where infinite scrolling, autoplay and push notifications were investigated as systematic manipulative features impairing free and informed choice. College students illustrate the developmental-vulnerability dimension: their unstable self-identity and incompletely matured executive functions render them susceptible to excessive use, associated with diminished academic concentration, weakened social skills, and heightened anxiety, depression and behavioural addiction. Longitudinal data \parencites[2]{8958037} specifically tracking recommendation-system effects on executive-function development in this cohort remain scarce, which constrains causal inference about platform-specific harm.

This report therefore defines algorithmic addiction as compulsive engagement emerging from the interaction of personalised recommendation logic, interface-level stimuli, and user vulnerabilities that limit reflective control. Social-media recommender research extends the concept beyond short video to broader feeds. Engagement-driven prioritisation \parencites[19]{3303310} amplifies misinformation and reinforces societal polarisation, so platforms simultaneously connect users and erode their cognitive autonomy. Collaborative-filtering architectures systematically amplify high-arousal material; Instagram's prioritisation of body-negative imagery and the associated decline in self-esteem among adolescent women of color illustrates harm generated while users remain only partially aware of the underlying curation.

Neurocognitive indicators are stark. Users exposed to hyper-personalised feeds show higher depressive-symptom incidence than those on chronological content. Neuroimaging links algorithmically amplified "doomscrolling" \parencites[3]{3303310} to reduced prefrontal cortex activity mirroring substance-dependence patterns. Compulsive-checking evidence in dopamine-driven feedback loops confirms that recommender logics and interface affordances repeatedly cue high-arousal responses that alter emotional regulation and impulse control over time \parencites[2]{3303310}[2]{8958037}. Hu's dataset recorded $23\%$ higher depressive-symptom incidence among users on hyper-personalised feeds compared with chronological controls \parencites[2]{3303310}[2]{8958037}.

Ethical and legal models supply normative boundaries. Kostenko's ethical-responsibility principle \parencites[61--62]{2977039} prohibits systems that manipulate behaviour or cause digital addiction, framing psycho-emotional comfort and informed decision-making as mandatory design objectives. The digital-neutrality principle requires documented decision logic and independent audit, recognising proven manipulation as a threat to digital sovereignty \parencites[109]{2977039}. Algorithmic addiction is thus treated as a violation of these duties: a state where optimisation for engagement and profit produces compulsive, autonomy-reducing usage \parencites[61--62]{2977039}[109]{2977039}[19]{3303310}[2]{8958037}.

Matrix factorisation-based recommenders frequently incorporate social ties to sharpen personalisation. Recent work demonstrates how both explicit friendships and inferred implicit connections can be fused into probabilistic matrix factorisation to improve predictive accuracy \parencites[2]{8026166}, a technique that can enhance relevance at the cost of reinforcing the same social-topology signals that drive engagement loops. No published study has tested whether social-tie fusion raises compulsive-use rates in the same way as purely content-driven signals.

\subsection{Positioning within Management and AI Ethics Scholarship}
\label{sec:positioning_mgmt_ai_ethics} Management scholarship on AI treats ethical governance as an organisational learning and control problem. Communities of Practice (CoP) are conceptualised as social learning systems where professionals develop capabilities through participation and joint problem-solving \parencites[3]{8200246}. Ethical AI in marketing is context-dependent, developing continually through inductive knowledge construction; fixed rulebooks cannot capture it. Reflexive thematic analysis is proposed to synthesise fragmented evidence on how organisational mechanisms support responsible AI implementation \parencites[8]{8200246}. The implication for this thesis is direct: an audit framework must be embedded in governance routines as an ongoing practice, never issued as a one-off checklist.

CSR is positioned as the governance bridge. It formalises CoP-generated insight into commitments, policies and accountability practices \parencites[16]{8200246}[11]{8200246}, converting moral awareness into routines for AI design, data management and performance measurement. Algorithm audits, AI ethics committees and CSR metrics make ethical aspirations enforceable inside firms \parencites[16]{8200246}. Stakeholder trust is the central outcome variable, resting on transparency, privacy, fairness and responsible personalisation; technical efficiency alone does not earn it. Consumer acceptance of AI-driven advertising \parencites[16]{8200246}[11]{8200246}[3]{8200246} is conditioned on perceived control and clarity, implying any recommender audit framework must address these dimensions to be managerially relevant.

Parallel AI governance proposals supply an external normative frame. Kostenko's AI Law Model gives concrete legal form to human dignity, autonomy and non-discrimination through mandatory procedures: notification of AI use, explainability, human control guarantees over critical decisions, and protection of vulnerable groups \parencites[9]{2977039}, backed by national AI expert councils, accredited auditors and dynamic risk reclassification. This resonates with management calls for independent audits and external oversight committees \parencites[9]{2977039}[9]{7461574}. Evidence on how well those councils and auditors perform under opaque black-box conditions comes mainly from financial and healthcare sectors \parencites[9]{2977039}; whether social media platforms would accept comparable oversight is untested.

Digital consumer and financial-ethics research anchors this further. AI-embedded financial decision support raises explainability and automation-bias concerns \parencites[3]{4245543}[2]{2855351}. Credit-decisioning evidence shows models trained on proxy variables can reproduce gender and intersectional discrimination even without protected attributes \parencites[9]{5442231}[7]{6059812}, meaning fairness audits focused only on overt attributes miss compounded subgroup disadvantage. Methodologically, ethical AI marketing and dark-pattern regulation studies settle on secondary qualitative designs synthesising regulatory texts and survey evidence where proprietary data access is constrained \parencites[6]{6295677}[8]{8200246}, echoed by constructivist grounded theory treating platforms as moral environments shaping identity and intention-behaviour gaps \parencites[3]{4245543}[5]{4245543}. These approaches validate this thesis's reliance on literature-based inquiry for theory-driven audit tools under black-box opacity. Personalisation-based dark-pattern research shows interface design induces decision distortion via loss aversion and anchoring \parencites[4]{7437483}, implying audit protocols must test manipulative elements and their cognitive mediators directly. Expert-AI-interaction studies show interfaces that surface uncertainty and let users reweight evaluative criteria preserve professional agency \parencites[6]{4268459}, a practical check that CSR-framed audits should formally adopt instead of leaving it to individual discretion.

\section{Literature Review and Theoretical Foundations}

\subsection{Architecture of Social Recommender Systems}

\subsubsection{Core Components of Recommendation Pipelines}
\label{sec:core_components_pipelines} Three interlocking layers determine what a recommendation pipeline surfaces: data ingestion, model computation and governance instrumentation. In social environments, data capture is anchored on implicit feedback: clicks, viewing time and browsing activity that are proxy indicators of interest without directly encoding satisfaction \parencites[21]{6249388}. Interaction data typically enters via distributed messaging into cloud data lakes. Historical bias enters at this stage: neutral variables like geolocation reconstruct protected attributes despite formal omission of sensitive fields \parencites[4]{8907572}, and social recommender pipelines relying on similar behavioural telemetry face the same structural risk.

Pre-processing transforms raw logs into model-ready vectors. Credit-scoring bias-mitigation engines operate in three phases: reweighting at pre-processing, adversarial debiasing during training, and reject-option classification at post-processing \parencites[4]{8907572}, aligning outputs with fairness constraints while preserving predictive performance. Core recommender models then operate on these signals. Content-based techniques retrieve items similar to prior interactions, limiting discovery. Collaborative filtering better supports serendipity but can surface niche items unevenly, with algorithms like Bayesian Personalised Ranking biased toward already-popular items.

Evaluation layers translate model behaviour into measurable outcomes. Beyond-accuracy metrics (diversity, novelty, catalogue coverage) assess whether a system broadens exposure. Bandit-based recommenders build exploration directly into the algorithm to counteract "rich get richer" feedback loops, though the correlation between these metrics and actual user perception remains unclear \parencites[15]{2190386}. Dataset curation shapes outcomes further. Benchmarks like MovieLens, Amazon and MIND each emphasise different properties: explicit ratings, implicit scale, or temporal and trust relationships. Public availability enables reproducible evaluation \parencites[21]{6249388} but introduces popularity and demographic biases that constrain generalisability to real-world social settings, a gap that field-deployment audits are better placed to close than benchmark comparisons alone.

Governance components (explainability, audit logging and bias mitigation) are increasingly embedded alongside these stages from the outset, as the logging architectures discussed below demonstrate \parencites[4]{8907572}.

\subsubsection{Feedback Loops between User Behavior and Ranking}
\label{sec:feedback_loops_user_ranking} Feedback dynamics arise whenever user interactions are recycled as training signals determining subsequent rankings. Short-video ecosystems illustrate this directly. Recommendation engines form a closed loop: "personalized $\rightarrow$ need satisfaction $\rightarrow$ behavioral data accumulation $\rightarrow$ algorithm optimization $\rightarrow$ re push." More precise targeting increases dependence while "information cocoons" narrow exposure, and engagement metrics function simultaneously as optimisation targets and as data sources entrenching the patterns they measure.

Neuropsychological mechanisms accentuate this. Dopamine release is tied to anticipation, surging before the reward arrives; the completed event itself produces less response, driving the "pursuit $\rightarrow$ transient satisfaction $\rightarrow$ further pursuit" cycle. Prolonged stimulation induces $\Delta$FosB expression \parencites[3]{8958037}, altering reward-circuit gene expression and weakening self-regulatory capacity especially among adolescents. Ranking policies that weight recent watch time do more than sequence content; they shape the neurocognitive conditions under which the next click occurs.

Broader social-media evidence connects such loops to documented outcomes. Engagement-oriented prioritisation \parencites[19]{3303310} amplifies misinformation and polarisation while users remain largely unaware of the mechanism. TikTok's system reinforces filter bubbles while improving loyalty. YouTube's recommender, responsible for over $70\%$ of views, prioritises watch time often at diversity's expense. Instagram and Facebook amplification of partisan material in election contexts intensified polarisation \parencites[2]{5142478}[23]{5142478}[29]{5142478}[19]{3303310}. Infinite scrolling and variable reward schedules correlate with $22\%$ higher compulsive-checking rates \parencites[2]{3303310}, and neuroimaging links prolonged engagement to reduced prefrontal activity resembling substance dependence, indicating a neurological substrate; the pattern is rooted in reward-circuit physiology.

Algorithmic bias compounds this. Systems trained on skewed interaction data reproduce cultural stereotypes. Regulatory analysis under the DSA/AI Act finds engagement-driven biases in core algorithms remain largely unaddressed despite transparency reporting \parencites[5]{5142478}[8]{5142478}[9]{5142478}, so feedback loops privileging sensationalist content persist. Proposed responses include fairness metrics and explanatory AI tools exposing why items appear in feeds.

The bounded-rationality research adds an interpretive layer: users satisfice and rely on availability and anchoring shortcuts \parencites[4]{3232181}, so click signals misrepresent stable preferences and instead reflect transient cognitive constraints. These heuristically biased traces are what ranking models learn from, and fairness audits should treat that distortion separately from downstream algorithmic output.

\subsubsection{Data Flows, Logging and Profiling Practices}
\label{sec:data_flows_logging_profiling} Consider AI-mediated educational systems as an illustrative case. IoT-enabled learning platforms aggregate biometric logs (voice, keystroke dynamics, facial recognition), behavioural analytics and network metadata into unified datasets distinguishing genuine activity from fraud \parencites[5]{0751607}. Pre-processing pipelines impute and normalise these into model-ready profiles. Financial AI architectures show a parallel high-volume, multi-tenant pattern: transaction streams ingested via Kafka into row-level-secured storage still admit historical bias at ingestion, since neutral variables like geolocation become precise demographic proxies within deep neural networks \parencites[4]{8907572}. A telemetry-parsing layer exposes incoming vectors to a domain-specific bias-mitigation engine so ethical evaluation runs inside the live flow; retroactive batch review becomes a fallback at most. Social recommender systems with similar ingestion architectures inherit the same bias risk.

Logging practices carry a distinct regulatory weight of their own. Financial systems run asynchronous SHAP/LIME computations in parallel with inference, hashing explanation vectors onto immutable ledgers for tamper-evident audit \parencites[4]{8907572}. HITL grading deployments time-stamp every AI and instructor decision with reviewer identifiers \parencites[16]{6805789}[17]{6805789}, making overrides traceable at the level of individual rubric dimensions. Ethically-grounded education frameworks extend logging toward explicit privacy guarantees: anonymisation, encrypted audit-log storage, opt-out options, and clear "AI generated" vs. "instructor reviewed" attribution aligned with FERPA and GDPR \parencites[20]{6805789}[5]{0751607}.

Regulatory AI law models generalise these into structured duties for high-risk systems. The TSAI regime requires owners to preserve event logs and comply with data-protection and transparency requirements, implemented through an algorithmic passport specifying data sources, legal grounds, log storage periods and audit information \parencites[242]{2977039}[244]{2977039}. For AI in public communications, auditable interaction logging distinguishes permitted dual-use tools from prohibited manipulative systems; logging intensity scales from increased to critical risk tiers, and platforms amplifying manipulative content without such safeguards bear public responsibility, with independent audit mandated to verify watermarking, labelling and log-based evidence collection \parencites[95]{2977039}[97]{2977039}. Privacy impact assessments and routine independent audits are practical complements to these cryptographic safeguards \parencites[6]{5988547}, treating logging as an operational layer running continuously alongside inference, with compliance records produced as a byproduct of that continuous monitoring.

\subsection{Engagement-Optimized Interface Design in Social Platforms}

\subsubsection{Infinite Scroll and Continuous Content Streams}
\label{sec:infinite_scroll_streams} Remove the stopping cue, and the user keeps scrolling. Infinite scroll achieves this by coupling a single low-effort gesture to a personalised next item, bypassing the chronological queue that once imposed a natural stopping point \parencites[2]{8958037}[3]{8958037}. Each scroll both delivers content and refines the model. The closed cycle, "personalized $\rightarrow$ user need satisfaction $\rightarrow$ behavioural data accumulation $\rightarrow$ algorithm optimization $\rightarrow$ re push," depends on this frictionlessness to remain active.

Neuroscience explains why frictionlessness is not neutral. Dopamine surges at the moment of expectancy before the next item, producing a "pursuit $\rightarrow$ transient satisfaction $\rightarrow$ further pursuit" cycle. Prolonged stimulation induces $\Delta$FosB expression \parencites[3]{8958037}[2]{8958037} in reward circuitry, undermining self-regulatory capacity, particularly among adolescents whose executive functions are still maturing. Variable reward schedules, identified \parencites[2]{3303310} as core elements of "dopamine driven feedback loops," yield $22\%$ higher compulsive-checking rates; neuroimaging links prolonged use to reduced grey-matter density in impulse-control regions, the same structural change observed in substance-dependence populations.

Regulatory analysis treats continuous streams as potentially manipulative design. The EU DSA investigated TikTok's infinite scrolling, autoplay and push notifications as systematic manipulative features impairing free decision-making. Article 25 DSA prohibits manipulative interfaces outright. Chinese algorithmic-recommendation rules separately require providers to offer non-personalised options and functions to disable recommendation entirely \parencites[2]{8958037}[7]{8958037}. TikTok's "For You Page" is described by DSA regulators \parencites[16]{6502223} as an exceptionally effective, opaque attention-capture system, notable precisely because users cannot identify the signals driving its curation.

Proposed mitigations target frictionlessness directly. Randomised content or customisable filters can disrupt information cocoons. Active exit nodes, usage-duration reminders and timed pauses reintroduce deliberation \parencites[7]{8958037}. "Why this content" explanations move from black-box manipulation toward visible, DSA-aligned assistance \parencites[7]{8958037}[8]{5142478}. This thesis treats continuous streams as ethically problematic specifically where infinite scroll combines engagement maximisation, ranking opacity and absent exit cues for vulnerable groups \parencites[2]{8958037}[3]{8958037}[2]{3303310}[16]{6502223}.

\subsubsection{Autoplay, Short-Form Clips and Sound Design}
\label{sec:autoplay_shortform_sound} Every completed clip immediately triggers another. That is the design intention of autoplay: deliberate stopping requires effort; continued consumption does not \parencites[2]{8958037}[3]{8958037}. Dominant formats are 15--60 second vertical clips chained through infinite swipe navigation. Douyin usage runs near 60 minutes across 21+ daily sessions, with $41.3\%$ of users opening the app multiple times daily and $56.9\%$ exceeding 60 minutes \parencites[3]{8958037}. These are not edge-case patterns; they describe the median user's daily experience.

For college students and adolescents with immature self-control, this pattern is associated with heightened addiction susceptibility, diminished academic concentration, and increased anxiety and depression \parencites[2]{8958037}[6]{8958037}. Sound design amplifies the effect: dynamic music paired with vertical clips produces a "sunflower seed effect" \parencites[3]{8958037} that disrupts default-mode-network functioning and compromises self-regulatory capacity. Neurobiologically, the content-switching mechanism resembles a slot machine. Each swipe carries unpredictable payoff. Dopamine surges at the moment of anticipation, before the clip begins. Chronic stimulation downregulates dopamine receptors and elevates reward thresholds \parencites[4]{8958037}, reducing the hedonic value of natural rewards in a self-reinforcing cycle.

Broader social-media evidence corroborates this: the same compulsive-checking and reduced-prefrontal-activity pattern documented above for infinite scroll \parencites[2]{3303310} recurs here, with autoplay and high-arousal soundscapes constituting a stimulus package tuned to neural addiction pathways; user preference does not drive the design, the design shapes it.

Regulatory sources treat this as potential manipulative design. The DSA investigated TikTok's infinite scrolling, autoplay and push notifications on precisely these grounds \parencites[2]{8958037}. Kostenko's AI Law Model prohibits gamification of addiction and exploitation of age-related vulnerabilities \parencites[223]{2977039}, mandating for minors secure defaults, bans on endless scrolling and autoplay, and visible "no personalization" switches. Reliability provisions require systematic incident logging and independent audits for such systems \parencites[57]{2977039}[79]{2977039}. Autoplay, short clips and sound design sit at the intersection of dopamine-driven feedback loops and explicit regulatory concern, especially where adolescents form the core user group \parencites[2]{8958037}[3]{8958037}[2]{3303310}[223]{2977039}.

\subsubsection{Notification Systems and Re-Engagement Triggers}
\label{sec:notification_reengagement_triggers} A notification can call a user back into the personalised cycle before they have decided to return. Push notifications and streak reminders operate alongside infinite scrolling and autoplay as core elements of an "algorithmic addiction loop" \parencites[2]{8958037}[3]{8958037}, acting as re-entry cues that simultaneously refresh the data used for further targeting.

Neurocognitively, the harm begins before the app opens. Notifications signalling potentially rewarding content trigger the "pursuit $\rightarrow$ transient satisfaction $\rightarrow$ further pursuit" cycle through anticipatory dopamine release. High-frequency stimulation induces $\Delta$FosB expression \parencites[3]{8958037}[2]{8958037} that supports addiction-like states, a particular concern for adolescents whose executive functions are still maturing. Compulsive-checking behaviour documented in broader "dopamine driven feedback loop" research \parencites[2]{3303310} necessarily interacts with notification policies that repeatedly invite re-entry.

Regulatory responses address this directly. Kostenko's AI Law Model prohibits gamification of addiction and exploitation of age-related vulnerabilities; for minors it bans endless scrolling, autoplay, streaks and forced quests outright, restricting notifications to minimal, neutral informational messages with visible "no personalization" switches \parencites[223]{2977039}. The confidentiality principle further prohibits using personal data for profiling aimed at influencing users without explicit consent \parencites[44]{2977039}, directly constraining notification pipelines built on fine-grained behavioural histories. Dark-pattern research confirms that nagging notifications wear down resistance and generate decision fatigue, exploiting System 1 heuristics over deliberative System 2 processing \parencites[2]{6295677}[8]{6295677}.

These findings point to the same working taxonomy. High-frequency, personalised, difficult-to-disable prompts are high-risk for vulnerable users. Low-frequency, informational, consent-aligned notifications are lower risk \parencites[2]{8958037}[3]{8958037}[2]{3303310}[223]{2977039}. Clinical research on adolescents adds a caution: many conceal severe distress and seek coping strategies without involving caregivers \parencites[9]{2750607}, so notification-based re-engagement aimed at minors risks aggravating hidden crises and needs explicit safeguards. Since the underlying architectures rely on offline clustering and sequence mining to surface highly contextual re-engagement moments \parencites[2]{7063165}, limiting sequential, context-aware personalisation in notification pipelines may matter as much as restricting message content itself.

\subsection{Behavioral Economics and Choice Architecture}

\subsubsection{Default Effects and Friction Costs in Digital Journeys}
\label{sec:default_effects_friction_costs} India's IRDAI banned pre-checked add-on insurance defaults in 2022 \parencites[8]{6295677} precisely because preselected options reduce cognitive effort enough to produce uninformed consent at population scale. Social platforms apply the same mechanism to engagement-maximising settings so that status-quo bias keeps users there without an affirmative choice \parencites[152]{1070605}[153]{1070605}. Roach-motel and subscription-trap patterns push exit costs up, and $47\%$ of users report cancellation difficulty \parencites[4]{6087241}[8]{6295677}; "nagging" notifications then create decision fatigue steering users back toward platform-preferred options. Dual-process cognition explains the synergy: defaults hold users in fast System 1 processing while sludge suppresses the reflective System 2 evaluation that meaningful consent decisions require \parencites[4]{6295677}. Counter-nudges, misleading button placement, and obfuscatory language can restore platform-aligned defaults without technically overriding regulatory requirements \parencites[9]{1070605}[12]{6295677}, which is why DMA research finds that enforcement intensity, independent of rule text, determines observed friction levels \parencites[152]{1070605}[160]{1070605}. Default bias introduced at the consent layer propagates downstream into content selection \parencites[3]{0123790}, and survey evidence shows users widely recognise algorithmic shaping yet still do not trust AI-curated content \parencites[20]{7646767}. Default settings and friction costs are therefore central manipulation-risk indicators; bright-line prohibitions and mandatory UX audits are the appropriate instrument, given how poorly case-by-case adjudication has contained them \parencites[12]{6295677}.

\subsubsection{Salience, Scarcity and Loss Aversion in Interface Design}
\label{sec:salience_scarcity_loss_aversion} Indian dark-pattern audits find "false urgency/scarcity" and "interface interference" on approximately $70\%$ of audited platforms \parencites[8]{6295677}. Prospect theory's loss asymmetry explains the mechanism: a countdown timer converts a routine purchase into a fabricated emergency by drawing System 1 attention and suppressing reflective appraisal \parencites[135]{8745065}[136]{8745065}. Executive interviews confirm the targeting logic is deliberate: budget-conscious consumers receive "pay more later" narratives while higher-income segments receive exclusivity framing \parencites[138]{8745065}, a vulnerability-targeting distinction unrelated to product relevance and described as routine competitive practice \parencites[135]{8745065}. The resulting architecture layers cognitive triggers, algorithmically-tailored timing, and compressed decision windows \parencites[139]{8745065}. DSA Recital 69 flags vulnerability-exploiting advertising as producing severe adverse effects; Article 26 mandates targeting disclosure; and sensitive-data profiling is categorically banned \parencites[18]{6941611}. These provisions draw the line between ordinary persuasion and the manipulation this taxonomy documents \parencites[12]{6295677}[390]{7896734}, and recommender personalisation shaping cultural-content distribution at scale \parencites[3]{5403359} sits squarely on the wrong side of it when urgency cues amplify the same cognitive vulnerabilities it surfaces.

\subsubsection{Temporal Discounting and Instant Gratification Online}
\label{sec:temporal_discounting_instant_gratification} Average Douyin usage runs at approximately 60 minutes per day across more than 21 daily sessions \parencites[3]{8958037}, a pattern micro-duration clips and infinite swipe navigation produce by making "just one more" nearly effortless. Micro-duration formats align structurally with temporal discounting: immediate low-cost rewards outcompete distant alternatives, producing measurable correlations with diminished academic concentration among adolescents \parencites[2]{8958037}[3]{8958037}. Dopamine surges before consumption; with chronic stimulation, reward thresholds rise and the hedonic value of slower learning-oriented activities steadily falls \parencites[4]{8958037}. Hu documents the cumulative toll: $22\%$ higher compulsive-checking rates and $23\%$ higher depressive-symptom incidence on hyper-personalised feeds \parencites[2]{3303310}[20]{3303310}, with neuroimaging linking doomscrolling to reduced prefrontal activity resembling substance-dependence profiles. A $2{,}239$-participant I-PACE study found past and present temporal focus positively correlated with addiction, and re-ranking approaches can moderate popularity biases while retaining acceptable accuracy \parencites[3]{9051044}, confirming that maximising instant engagement is a deliberate design choice with measurable neurological consequences. DSA Article 25 bans interfaces that "substantially impair" free decisions \parencites[20]{6941611}[2]{8958037}, and literacy surveys confirm susceptibility persists despite high self-reported digital confidence \parencites[31]{9360054}. Platform-side structural change \parencites[99]{2977039} addresses the source; individual countermeasures do not.

\subsection{Nudging, Dark Patterns and Ethical Boundaries}

\subsubsection{From Helpful Nudges to Manipulative Dark Patterns}
\label{sec:helpful_nudges_to_dark_patterns} Ethical choice architecture and dark-pattern manipulation share structural form but differ in what they do to user comprehension and reversibility \parencites[32--33]{6087241}. A recommender that hides cheaper alternatives or makes refusal deliberately difficult defeats a choice; it does not clarify one. Drip pricing and confirm-shaming trigger fast System 1 processing; obstruction and forced action constitute "sludge" \parencites[2]{6295677}[4]{6295677} that suppresses the System 2 evaluation protective decisions require. Large-scale marketing crawls document the scale: asymmetric cookie consent and hidden fees increase immediate compliance while eroding long-term customer equity \parencites[33]{6087241}. Indian evidence is direct: $87\%$ of major platforms deploy at least one dark pattern, yet only $39\%$ of users recognise the term while $68\%$ report experiencing deception \parencites[2]{6295677}. India's Advertising Standards Council has defined thirteen dark-pattern categories \parencites[3]{6295677}, and judicial rulings against concealed BookMyShow fees confirm interface-level deception falls under consumer law. Platforms can neutralise regulatory choice screens with misleading button placement and obfuscatory language, preserving nominal autonomy while undermining its practical exercise \parencites[159]{1070605}. For the Ethical AI Audit Toolkit in Section 5, interface elements qualify as acceptable nudges when they simplify decisions without degrading comprehension or reversibility, and fall into dark-pattern categories when they rely on information distortion, interface asymmetry, social or temporal pressure, choice obstruction, or forced action \parencites[32--33]{6087241}[2]{6295677}[4]{6295677}.

\subsubsection{Common Dark Interface Strategies in Social Media}
\label{sec:common_dark_interface_social} Interface interference and obstruction appear on approximately $70\%$ and $63\%$ of major Indian platforms respectively; drip pricing, bait-and-switch, and privacy zuckering affect $75\%$, $48\%$, and $44\%$ of consumers in the same content analysis \parencites[7]{6295677}[9]{6295677}. Social platforms adapt these patterns from e-commerce and embed them in feeds, consent flows, and notifications coupled to recommender engines \parencites[7]{6295677}[3]{6087241}. Even where regulators introduce choice screens or interoperability mandates, dominant platforms neutralise them with friction and defaults \parencites[3]{1070605}[10]{1070605}. AI-enabled environments add escalation patterns: "adaptive UI and friction shifting" uses reinforcement learning to trap users in roach-motel cancellation journeys, while "chatbot coercion" deploys conversational agents to frame cancellation as socially deviant \parencites[11]{7896734}. Large language models mass-produce indistinguishable fake testimonials; confirm-shaming personalises guilt-inducing opt-out language; fabricated "friends liked this" signals manufacture social pressure at scale \parencites[11]{7896734}. The self-declared compliance figures are striking: $81\%$ of platforms claimed to be "dark pattern free" during CCPA audits while $87\%$ deployed at least one pattern, and complaint resolution reached only $28\%$ \parencites[7]{6295677}[12]{6295677}. Interface interference, obstruction, AI-driven friction shifting, and conversational coercion each directly impair autonomous exit and are classified as central items in this report's working taxonomy \parencites[3]{6087241}[7]{6295677}[11]{7896734}.

\subsubsection{Moral and Legal Contours of Manipulative Design}
\label{sec:moral_legal_manipulation} Kostenko's ethical-responsibility principle requires AI systems to counter manipulative design through mandatory notification, explainability, and human control over critical decisions \parencites[8]{2977039}[61--62]{2977039}, with addiction-causing systems categorically prohibited. The TSAI regime provides the legal structure: infrastructure access is conditioned on registration, risk classification, and regular algorithmic-passport audits, with violations triggering fines or identifier suspension \parencites[243]{2977039}[244]{2977039}. Indian e-commerce evidence confirms the patterns are systemic: interface interference, obstruction, and drip pricing exploit dual-process cognition at scale, prompting calls for bright-line prohibitions and mandatory UX audits \parencites[2]{6295677}[4]{6295677}[13]{6295677}. The DSA bans interfaces that "substantially impair" free decisions; the DMA bans gatekeeper interface pressure; and the anticipated Digital Fairness Act targets addictive product design directly \parencites[389]{7896734}[390]{7896734}. Georgia's "Maxbosch"/BOSCH case sanctioned false countdown timers without AI-specific statutes \parencites[390]{7896734}[391]{7896734}, showing existing consumer protection law already reaches manipulative interfaces. For this report's audit framework: ethical design is documentable through algorithmic passports and audit trails; manipulative design, defined as intentional exploitation of cognitive vulnerability with opaque targeting and addiction-inducing engagement, falls within what AI law models seek to prohibit outright or subject to registration and sanction \parencites[8]{2977039}[61--62]{2977039}[389]{7896734}[13]{6295677}.

\subsection{Motivational Psychology and Self-Determination}

\subsubsection{Autonomy, Competence and Relatedness in Digital Contexts}
\label{sec:autonomy_competence_relatedness_digital} Kostenko's ethical-responsibility principle establishes that meaningful autonomy requires more than an off-switch: design must preserve users' capacity to understand and contest algorithmic influence \parencites[61--62]{2977039}, and systems causing addiction-inducing manipulation are prohibited outright. The accompanying right to a "psycho-emotionally safe digital environment" includes audit tools for one's own activity \parencites[119]{2977039}[2]{8958037}, placing infinite scroll and high-frequency notifications within the scope of autonomy harm when they undermine critical thinking. Neural-network recommendation combined with variable reward schedules pushes users into "auto pilot" scrolling, already flagged by the DSA as potential manipulative design, with disproportionate impact on college students whose executive function is still maturing \parencites[2]{8958037}. Neuroimaging evidence confirms that autonomy constraints operate at the neurocognitive level, extending well below the consent-screen layer \parencites[2]{3303310}. Competence faces a parallel gap. Literacy surveys report high self-assessed confidence in routine digital tasks but a mean score of only $2.31$ out of $5$ for recognising manipulation \parencites[31]{9360054}. Kostenko's digital-hygiene rights address this "second-level divide" directly \parencites[119]{2977039}, and the Algorithmic Impact Navigator's broad stakeholder-engagement model confirms that recommender audits should involve users and civil society alongside technical teams \parencites[2]{2742041}[25]{4900172}. Relatedness manipulation is the least visible vector. Inflated engagement metrics via bots on Douyin and Xiaohongshu, combined with undisclosed influencer sponsorship, convert social validation into a tool for manufacturing artificial consensus \parencites[13]{7160486}, and popularity-amplifying recommendation logic embeds those false signals directly into the ranked feed. Annual external audits and accessible complaint mechanisms, backed by revocable-consent requirements, are the operational remedies the literature proposes \parencites[119]{2977039}[81]{2977039}[7]{8453044}. User-awareness campaigns reduce susceptibility at the margins; the structural conditions degrading autonomy, competence, and relatedness remain intact until the design architecture itself changes \parencites[61--62]{2977039}[119]{2977039}[31]{9360054}[13]{7160486}.

\subsubsection{Impact of Algorithmic Environments on Intrinsic Motivation}
\label{sec:impact_intrinsic_motivation} Blooket's token system shows the design choice concisely: tokens function as a "currency for creative agency" when embedded in goal-oriented narratives, satisfying autonomy and competence; poorly calibrated extrinsic rewards produce the opposite \parencites[7]{9283805}, crowding out internal satisfaction through overjustification. Short-video recommendation ecosystems represent the corrosive configuration at scale. Dopamine surges at the moment of scrolling anticipation drive a "pursuit $\rightarrow$ transient satisfaction $\rightarrow$ further pursuit" cycle concentrated among adolescents with lower self-control \parencites[2]{8958037}[3]{8958037}. Hu documents the downstream effect: $22\%$ higher compulsive checking and reduced prefrontal activity in impulse-control regions \parencites[2]{3303310}, making initiation of demanding future-oriented tasks harder even when underlying values are intact. Kostenko's ethical-responsibility principle bans addiction-causing systems \parencites[61--62]{2977039}; the confidentiality principle prohibits profiling-driven influence without explicit consent \parencites[44]{2977039}. The anticipated Digital Fairness Act's addictive-design focus \parencites[390]{7896734} and dual-process dark-pattern research on System 1/System 2 exploitation \parencites[6]{6295677} converge on the same harm: architectures that substitute platform-defined reward schedules for users' own goals, rarely disclosed as such. Gradually fading extrinsic-reward prominence internalises habits, shifting motivation away from point dependency \parencites[7]{9283805}, and an AI-augmented mental-health intervention in Nigeria confirmed that human intermediaries remain essential to preserve self-directed coping \parencites[9]{5419804}.

\subsubsection{Experiences of Control, Flow and Compulsion in App Use}
\label{sec:experiences_control_flow_compulsion} $1.074$ billion Chinese short-video users, representing $95.4\%$ of netizens, show excessive use linked to diminished academic concentration and behavioural addiction, concentrated among adolescents with lower self-control \parencites[2]{8958037}. Micro-duration clips coupled with infinite swipe navigation create a closed delivery cycle drawing users into "auto pilot" scrolling \parencites[2]{8958037}[3]{8958037}. Dopamine surges during anticipation; the act of consumption yields only transient relief; chronic stimulation induces $\Delta$FosB expression, elevates reward thresholds \parencites[3]{8958037}[4]{8958037}, and weakens self-governed choice even during sessions that feel like effortless immersion. A $2{,}239$-participant I-PACE study found past and present temporal focus positively correlated with addiction while future orientation was protective \parencites[4]{8958037}, mediated by a "resistance backfire effect" under inhibitory self-control versus substitute-behaviour initiation under initiatory self-control. Broader social-media evidence records $22\%$ higher compulsive-checking rates and reduced grey-matter density in impulse-control regions \parencites[2]{3303310}[20]{3303310}. Kostenko's model prohibits algorithms causing psychological pressure or forming addiction \parencites[61--62]{2977039}[119]{2977039}, grounding the ban in an inalienable right to a digital environment free from exhaustion and manipulative influence. Audit practice should therefore classify subjective reports of "getting lost" in feeds and post-use exhaustion as misalignment signals requiring investigation.

\subsection{European Digital Governance Frameworks}

\subsubsection{Digital Services Regulation of Online Platforms}
\label{sec:digital_services_regulation_online_platforms} The Digital Services Act (DSA) reframes recommender systems as objects of systemic accountability. Platforms exceeding 45 million EU users ("very large online platforms," VLOPs) must conduct systemic risk assessments, submit to annual independent audits, offer users at least one non-profiling recommendation option, and disclose the main parameters driving recommendations \parencites[7]{1456534}[12]{6502223}. Three audit tiers emerge: internal risk assessments by providers, independent compliance audits under Chapter III, and vetted-researcher data access under Article 40(4) \parencites[18]{3195574}[19]{3195574}.

Whether these obligations have teeth is contested. Audits risk becoming procedural formalities where regulators \parencites[9]{7461574}[9]{1456534} lack the technical expertise to interrogate ranking algorithms, a fragmentation already visible in uneven GDPR enforcement across member states. Analysts nonetheless anticipate a "Brussels Effect," with DSA transparency and audit norms establishing a de facto global benchmark \parencites[8]{1456534}[7]{5162643}. The systemic risk assessment requirement, non-profiling option, and audit-tier structure form the compliance floor this report's Ethical AI Audit Toolkit must internalise \parencites[12]{6502223}[2]{1456534}.

\subsubsection{Artificial Intelligence Risk Governance in the EU}
\label{sec:ai_risk_governance_eu} Four risk tiers define the EU AI Act's architecture. Unacceptable risk (social scoring, real-time biometric surveillance, exploitation of psychological vulnerabilities) is banned outright. High-risk systems in education, employment and law enforcement require conformity assessment, technical documentation and human oversight. Limited risk carries transparency duties; minimal risk faces no mandatory obligations \parencites[10]{2977039}[12]{6502223}. Kostenko characterises this EU regime as the most systemic globally, contrasting it with the United States' sectoral soft-law approach and China's centralised, state-centric model. ISO/IEC JTC 1/SC 42 standards, for example ISO/IEC 24027 on bias and 24028 on trustworthiness, converge conceptually with the Act's risk categories but remain voluntary \parencites[4]{8589961}[5]{8589961}.

\begin{table}[H] \centering \caption{Selected AI risk governance features by region} \begin{tabular}{lccc} \hline Feature & EU AI Act & US federal & China \\ \hline Legal form & Binding regulation & Soft guidance & Administrative regs \\ Risk tiers & 4 level scheme & None unified & Domain specific \\ High risk duties & Conformity, oversight & Sectoral & State centric \\ Audit capacity & Mandated, still weak & Fragmented & Centralised \\ \hline \end{tabular} \end{table}

Institutional capacity remains the critical bottleneck. Auditing opaque, black-box models is still weak even within the EU \parencites[9]{7461574}[2]{1456534}, and national competent authorities face real pressure to develop algorithmic-systems expertise within the Act's compliance deadlines. The AI Act interlocks with the DSA on systemic risk duties for VLOPs and with the DMA on gatekeeper competition duties, pushing algorithmic accountability toward infrastructural questions of market power alongside, and at times displacing, individual rights.

Concrete harm is documented in adjacent domains. Algorithms trained on skewed school data have systematically underrated lower-income students, with exposed labels measurably reducing teacher engagement \parencites[2]{6362702}, a feedback loop between training-data composition and real-world disparity that the Act's dataset and conformity assessment obligations are designed to interrupt.

\subsubsection{Policy Trajectory toward a Digital Fairness Regulation}
\label{sec:policy_trajectory_digital_fairness} A gap sits at the centre of current EU digital regulation: neither the GDPR, the DSA, nor the AI Act currently classifies social recommender algorithms as high risk in their own right \parencites[2248]{6502223}[8]{1456534}. Sanikidze describes a planned Digital Fairness Act (DFA), proposed for late 2026, that would regulate addictive product designs and unfair personalisation "across the board" \parencites[390]{7896734}, filling the gap between the DSA and the Unfair Commercial Practices Directive. If enacted at its proposed scope, the DFA would be the first instrument to treat interface-level addiction risk as a primary regulatory target in its own right, independent of privacy or competition enforcement.

Enforcement practice already previews what a fairness-focused regime might catch. Georgia's Competition and Consumer Agency has sanctioned dark patterns (false countdown timers, fabricated scarcity, misleading brand associations) under general unfair-practices law, though those cases have not yet reached AI-driven dark patterns specifically \parencites[390]{7896734}. Recent DSA and UCPD enforcement against platform X and Temu, including fines for deceptive verification flows, signals regulatory impatience with interface pressure of exactly this kind.

Provider-commissioned reviews vary widely in independence \parencites[13]{2574631}, meaning credible compliance will require evidence of substantive, replicable testing; declarative documentation will not suffice. For this thesis, which treats the DFA as an anticipated policy trajectory whose final statutory text remains unpublished \parencites[2248]{6502223}[390]{7896734}, the operational implication is to build audit dimensions for addictive design and unfair personalisation ahead of the DFA's formal text. Doing so ahead of formal enactment carries the acknowledged risk that the DFA's final scope differs from current proposals.

\subsection{Transparency, Data Access and Audit Obligations}

\subsubsection{Platform Transparency Reporting Requirements}
\label{sec:platform_transparency_reporting} Transparency reporting under the DSA is the primary control surface for recommender accountability. VLOPs must publish recurring reports on disputes, suspensions and moderator staffing under Article 24, disclose the main parameters driving recommendations under Article 27, maintain an advertising repository under Article 39, and grant data access to regulators under Article 40 \parencites[19]{3195574}[7]{1456534}.

\begin{table}[H] \centering \caption{Key DSA transparency reporting elements and audit implications} \label{tab:dsa_transparency_dims} \begin{tabular}{p{3.1cm}p{4.2cm}p{5.0cm}} \hline Regulatory hook & Required disclosure & Relevance for Ethical AI Audit \\ \hline Art. 24 reports & Disputes, suspensions & Indicator of downstream harms and appeal pressure linked to recommender outputs \\ Art. 27 parameters & Main recommendation signals & Entry point for interrogating engagement optimisation and profiling dependence \\ Art. 39 ad repository & Ad content and targeting info & Enables cross checking of content vs ad recommendation consistency \\ Art. 40 data access & Data sharing to regulators & Conditions for external, scenario based audits of systemic risks \\ \hline \end{tabular} \end{table}

Me{\ss}mer and Degeling read these clauses as the backbone of a multi-actor audit regime, but warn that without detailed implementing acts platforms could perform partial self-assessment while remaining formally compliant \parencites[6]{3195574}[19]{3195574}. Sebastian anticipates algorithmic impact assessments for high-risk domains resembling environmental impact reviews, with ongoing disclosure at each update cycle, because model drift can introduce new discrimination over time \parencites[15]{5442231}. Kostenko's AI Law Model generalises this into a transparency principle \parencites[70]{2977039} requiring disclosure of data sources, architectures, error probabilities and uncertainty scenarios, with covert or manipulative AI use treated as categorically unacceptable. All three converge on the same requirement: disclosure must be detailed enough that a second auditor could reproduce the test.

An internal audit dashboard should track recommendation exposure by content category, repeat exposure to flagged content, and non-profiled-option uptake. Translating Article 27 parameters and Kostenko's disclosure requirements \parencites[12]{6502223}[70]{2977039} into plain-language external statements lets regulators reconstruct how engagement objectives interact with fairness obligations.

\subsubsection{Data Access for Regulators and Independent Researchers}
\label{sec:data_access_regulators_researchers} Without data access, systemic auditing of recommenders is impossible. Article 40 of the DSA obliges VLOPs to grant access to Digital Services Coordinators, the Commission, and vetted researchers studying systemic risks such as civic-discourse effects \parencites[19]{3195574}[12]{6502223}. Me{\ss}mer and Degeling turn this into a scenario-based model, mapping specific risk narratives (eating-disorder promotion, for instance) to concrete, scenario-grounded Article 40 data requests \parencites[15]{3195574}. Kostenko's AI Law Model pushes the principle beyond EU platform law into cross-border "extraterritorial digital jurisdiction" \parencites[249]{2977039}: any AI system affecting a State's territory or market falls under that State's incident-reporting, audit-cooperation and evidence-preservation duties, transmitted through supply chains via flowdown clauses. In practice, this means a platform operating solely from outside the EU could still be audited on the basis of users it reaches within a member state's jurisdiction.

\begin{table}[H] \centering \caption{Data access regimes for AI and platform oversight} \begin{tabular}{p{3cm}p{5cm}p{5cm}} \hline Dimension & DSA / EU platforms & AI Law Model (Kostenko) \\ \hline Primary addressee & VLOPs with $45+$ million users & All AI systems with impact on State territory \\ Legal hook & Art. 40 data access, vetted researchers & Extraterritorial digital jurisdiction, cross border liability \\ Data objects & Platform data for systemic risk research; ads repository (Art. 39) & Technical materials, logs, audit results, algorithmic passports \\ Research access & Vetted researchers for systemic risk studies & API access to registry attributes for scientific research \\ Logging requirement & Implied through audit and risk assessment duties & Explicit event logging, evidentiary preservation, log based arbitration \\ \hline \end{tabular} \end{table}

Formal access rights do not execute themselves. Regulators still struggle with black-box models \parencites[9]{7461574}, and Me{\ss}mer and Degeling anticipate a fragmented mix of mandatory and voluntary audits, with no unified enforcement regime emerging yet \parencites[19]{3195574}. Work on AI-enhanced auditing in finance points toward human-AI collaboration, with AI handling large-scale analytics while human experts retain contextual judgement \parencites[3]{6539584}[6]{6539584}. That hybrid demands carefully designed interfaces and cognitive-load management for auditors, and without sustained investment in auditor capacity the formal Article 40 architecture will remain largely symbolic.

\subsubsection{Documentation, Logging and Traceability Expectations}
\label{sec:doc_logging_traceability} Kostenko's AI Law Model frames traceability as a foundational transparency and explainability requirement: any AI system affecting rights, security or wellbeing must be identifiable, traceable, explainable and accountable as a condition of lawful use, with decision logs, data pipelines and model parameters remaining interpretable enough for independent auditors to identify errors or bias \parencites[35]{2977039}[36]{2977039}.

Child-affecting systems raise this bar further, requiring a complete technical, methodological, legal and ethical dossier in which decision logs record configuration changes and human interventions; absent safety mechanisms or manipulative interfaces in child-facing contexts count as governance violations justifying suspension \parencites[214]{2977039}. Territorial sovereignty clauses extend traceability into a jurisdictional requirement \parencites[124]{2977039}: only AI systems registered in a National Register of AI Systems may operate within a State's territory, embedding logging and documentation directly into cross-border control. An unregistered system has no lawful path to deployment, whatever its technical merits.

Financial-AI systems stream SHAP-based feature attributions to a write-once immutable ledger, with circuit breakers switching models to human-supervised mode when drift crosses a threshold \parencites[7]{8907572}. HITL grading in education keeps audit trails of AI outputs and human overrides available for accrediting bodies, their evidentiary value dependent on tamper-evidence and retention period \parencites[15]{6805789}.

Me{\ss}mer and Degeling caution that without firmer implementing acts \parencites[6]{3195574}[19]{3195574}, platforms may default to minimal documentation that meets the letter but not the substance of the requirement. For an Ethical AI Audit Toolkit under an anticipated DFA the operational requirement set is: structured technical dossiers and algorithmic passports, tamper-evident logs, documented human-oversight and appeal channels, and independent access for regulators and vetted researchers \parencites[9]{7461574}[12]{6502223}. Continuous audit architectures that stream live, queryable evidence to oversight bodies can also improve governance itself \parencites[2]{4639902}[10]{7504644}, though confidentiality and false-positive concerns in that architecture remain unresolved.

\subsection{Child and Youth Protection in Algorithmic Environments}

\subsubsection{Age-Sensitive Risk Assessment and Design Standards}
\label{sec:age_sensitive_risk_design} Child-directed AI requires a qualitatively different assessment standard. Kostenko's AI Law Model mandates "secure by default" architecture: algorithmic logic, data categories and training processes must be understandable, verifiable and controllable across the system's entire lifecycle from design through deployment. This runs through a multi-stakeholder documentation set covering Child Impact Assessments, Child Safety Data Sheets, audit reports, and a minimum public disclosure baseline \parencites[214]{2977039}. Explainability is age-calibrated, requiring plain-language explanations at technical, functional and ethical levels, with children's rights tools such as profiling refusal and human review genuinely accessible and demonstrable in practice.

Absent safety mechanisms or manipulative interfaces in child-affecting contexts are classified as governance violations \parencites[223]{2977039}, empowering regulators to suspend systems and revoke certification. Age-appropriate design principles add concrete UX requirements: profiling, geolocation and biometrics disabled by default; bans on endless scrolling, autoplay, streaks and exit penalties for minors; capped notifications; and a visible "no personalisation" switch and one-click "why am I seeing this" explanation. Table~\ref{tab:child_risk_design} converts these into this report's working audit grid.

\begin{table}[H] \centering \caption{Working audit grid: age sensitive risk assessment and design standards} \label{tab:child_risk_design} \begin{tabular}{p{3.2cm}p{4.4cm}p{4.4cm}} \hline Dimension & Required evidence & Illustrative audit question \\ \hline Child Impact Assessment & Documented risk analysis, threat models, Child Safety Data Sheet & Is there a current Child Impact Assessment covering profiling, attention capture, and recommendation logic \parencites[214]{2977039}? \\ Secure by default & Architecture, data maps, disabled profiling/geolocation at default & Do defaults disable profiling, geolocation, biometrics, and addictive mechanics for minors \parencites[223]{2977039}? \\ Explainability & Multi level, age appropriate explanations & Can a child access "brief/detailed/technical" explanations of "why am I seeing this" in simple language \parencites[214]{2977039}[223]{2977039}? \\ Controls & No personalisation switch, parental controls, time limits & Are "lesson/sleep" modes and refusal of profiling available without service degradation \parencites[223]{2977039}? \\ Logging & Decision and event logs with minimum retention & Are tamper evident logs available to regulators and families in case of disputed decisions \parencites[214]{2977039}[191]{2977039}? \\ Governance & Internal and external child safety audits, regulator dossier & Has an external audit of child safety been completed and provided to authorities \parencites[214]{2977039}? \\ \hline \end{tabular} \end{table}

\subsubsection{Restrictions on Profiling, Targeting and Monetization}
\label{sec:restrictions_profiling_targeting_monetization} Kostenko's AI Law Model categorically bans targeted advertising personalised on a child's psychotype, neurobehavioural signs or predicted emotional state \parencites[20]{2977039}[224]{2977039}. Automated emotion and biometric tracking requires dual (child plus guardian) informed consent. Predictive profiling of a child's "future trajectories" is barred without qualification. Operators must offer "child information invisibility" modes, meaning service use without profiling backed by minimal, non-identifying logs, because children's consent is treated as fragile regardless of its formal collection method \parencites[223]{2977039}[44]{2977039}.

DSA Article 28 prohibits profile-based behavioural advertising targeting at minors once a platform can reasonably infer a recipient is a child \parencites[19]{6941611}. AI-powered health services for teenagers are barred from cross-platform tracking and third-party SDKs entirely \parencites[223]{2977039}, and monetised "digital twin" simulacra require revocable double consent; political or commercial use of minors' simulacra is prohibited outright \parencites[140]{2977039}. Regulators can suspend operation and impose liability on operators, suppliers and manufacturers simultaneously \parencites[223]{2977039}[44]{2977039}.

Any child-facing module combining profiling, targeting and monetisation is presumptively classified "prohibited or critical risk" \parencites[224]{2977039}[19]{6941611} unless the operator demonstrates information-invisibility modes and strict separation of minor data from advertising and recommendation pipelines, consistent with how the Kostenko model assigns the burden of proof to fragile-consent populations generally.

\subsubsection{Safeguards for Youth-Focused Platforms}
\label{sec:safeguards_youth_focused_platforms} Kostenko's model applies age-segmented standards. Early childhood (0--6) requires zero data collection by default, no personalisation or advertising, and local-only processing. Junior schoolchildren (7--12) face bans on loot-box monetisation and manipulative retention mechanics \parencites[217]{2977039}, including endless scroll, streaks and forced quests; AI decisions affecting this group must be explained in child-appropriate language, since retention mechanics causing acceptable friction for adults can produce disproportionate distress in 7-to-12-year-olds.

For ages 13--15 on teen social networks, feeds must default to "detox mode": no endless scroll, no autoplay, no aggressive return triggers. Notifications are capped at one per day; profiles are closed by default; a public Child Safety Data Sheet, quarterly Child Impact Assessments and annual independent audits are mandatory \parencites[221]{2977039}. Older teens (16--17) gain expanded self-determination, but predictive profiling and data monetisation remain prohibited. A further layer extends minimal-data-by-default and accessibility requirements to children with disabilities, displacement or economic hardship \parencites[219]{2977039}, since these groups face compounded exposure risk when algorithmic engagement combines with social or economic vulnerability.

Any youth platform combining profiling, behavioural advertising and addictive mechanics is classified "high or unacceptable risk" regardless of engagement metrics. A nationwide Register of Responsible Subjects of AI Systems backs this with accountability infrastructure, placing non-compliant operators on a Register of Violators \parencites[82]{2977039}. Applied evidence adds a caution: mechanics such as streaks that appear benign for adult users can edge into coercive nudging for minors \parencites[5]{4426226}, and clear anonymity controls measurably increase users' willingness to engage honestly with safeguards \parencites[5]{4426226}.

\subsection{Human-Centered Design and Autonomy Preservation}

\subsubsection{Respect for Agency and Meaningful Choice}
\label{sec:respect_agency_choice} Kostenko's ethical-responsibility principle defines the operative standard: design must keep "autonomy... [and] the right to an informed decision" intact, explicitly prohibiting systems that manipulate behaviour or suppress autonomy \parencites[61--62]{2977039}. A complementary digital-ecology right requires a "fair, transparent, humane" environment free from intrusive information flows \parencites[119]{2977039}. Together these frame autonomy as a positive design obligation: every feature must be testable against whether it expands or constrains the user's capacity for uncoerced choice.

Liu and Keshavarz define dark commercial patterns as designs that "steer, deceive, pressure, or obstruct consumers" in ways that impair "informed, preference consistent choice" \parencites[3]{6087241}[39]{6087241}. This thesis adopts that functional test as its agency-risk baseline. The DSA enforces the same boundary by requiring a non-profiled feed option and prohibiting interfaces that "substantially impair" free decisions \parencites[7]{1456534}[18]{6941611}. The anticipated Digital Fairness Act extends this toward addictive design and unfair personalisation \parencites[390]{7896734}, with OECD/ICPEN analysis framing manipulative design as a systemic autonomy threat \parencites[391]{7896734} that cannot be adequately addressed by individual consent mechanisms alone.

Literacy surveys complicate enforcement. Respondents show high self-assessed digital competence alongside a mean manipulation-detection score of only $2.31$/$5$, despite moderate trust in EU regulation \parencites[31]{9360054}. Liu and Keshavarz's "conversion-trust" framework translates this gap into three audit checks: information clarity, choice symmetry, and proportional reversibility \parencites[39]{6087241}. A European Commission study found $97\%$ of $45$ websites and $30$ apps using deceptive hierarchy or concealment practices, social media platforms included \parencites[13]{6941611}. Large-scale crawls independently confirm that cookie banners are asymmetrically designed: accepting tracking is systematically faster and easier than refusing it \parencites[33]{6087241}, a pattern consistent across multiple independently sampled datasets.

The AI Law Model closes what interface design leaves open, mandating independent ethical reviews, internal audit mechanisms and annual AI Ethics Impact Reports \parencites[61--62]{2977039}, plus bans on self-replicating addictive environments and suspension powers over non-auditable systems \parencites[119]{2977039}[108]{2977039}. Agency is treated as respected in this thesis only where non-profiled feeds are genuinely accessible \parencites[390]{7896734}, consent flows meet symmetry and reversibility criteria \parencites[33]{6087241}, and documented ethical audits address addiction risk \parencites[61--62]{2977039}[119]{2977039}.

\subsubsection{Explainability and Comprehensibility for Lay Users}
\label{sec:explainability_lay_users} Kostenko's transparency principle requires architecture, data sources, and implications to be "understood, accessible, verifiable, and explanatory" throughout the lifecycle, with covert AI use in socially significant contexts explicitly prohibited \parencites[67]{2977039}. The standard is stakeholder-specific: developers document internally, users need literacy-adapted interfaces, and regulators require a complete machine-readable dossier \parencites[68]{2977039}.

Kostenko's digital safety-of-the-child provisions make the requirement direct for minors: systems must deliver "meaningful, structured and age-appropriate" interpretations free of jargon, backed by decision logs and independent child-safety audits \parencites[214]{2977039}[223]{2977039}. HITL grading shows the toolkit working at scale: SHAP explanations highlighted score-contributing features, instructors confirmed or overrode with dual-logged actions, measurably increasing consistency and student trust \parencites[16]{6805789}[21]{6805789}. Sebastian's financial-fairness research adds SHAP-based adverse-action notices backed by ongoing algorithmic impact assessments \parencites[6]{5442231}[15]{5442231}, while Gerr proposes XAI "why this content" models for both user understanding and regulator-side distributional fairness checks \parencites[8]{5142478}[30]{5142478}.

Explanation quality matters independently of availability. Opaque logic triggers guilt, anger, or "ethical disengagement" \parencites[8]{4245543}[191]{4245543}: users stop attending to warnings altogether when they cannot interpret the reasoning behind them. Four audit expectations follow: stakeholder-differentiated explanation layers \parencites[68]{2977039}[214]{2977039}; multi-level explanations combining local reasons with global ranking descriptions \parencites[68]{2977039}[8]{5142478}; tamper-evident logging enabling contestation \parencites[214]{2977039}[241]{2977039}; and HITL mechanisms intelligible to both professionals and affected users \parencites[16]{6805789}[21]{6805789}. Journalism-ethics research frames this as institutional duty: verifiable AI-content labelling and audit logs are prerequisites for public confidence; their absence visibly undermines legitimacy \parencites[94]{2977039}[4]{7076960}.

\subsubsection{Meaningful Human Oversight of Automated Decisions}
\label{sec:meaningful_human_oversight} Kostenko's AI Law Model establishes the baseline: all lifecycle actors must implement internal ethical audit mechanisms and human-rights impact monitoring, with regulators empowered to suspend operation and annual AI Ethics Impact Reports mandated \parencites[61--62]{2977039}. HITL grading supplies a concrete implementation template: a four-layer framework lets instructors approve, adjust, or override AI-generated scores, with all changes dual-logged against the model's initial assessment \parencites[16]{6805789}. That traceability measurably strengthened trust, with $61\%$ of surveyed students reporting higher confidence because the decision path was transparent and contestable \parencites[8]{6805789}[20]{6805789}.

Business-analytics research reinforces the same structure. Human-in-the-loop, on-the-loop, and over-the-loop arrangements covering pre-implementation review, monitoring with intervention authority, and ex-post audits are prerequisites for AI Act/GDPR compliance against discriminatory lending or biased recruitment \parencites[7]{7461574}[9]{7461574}, echoed in technical governance frameworks requiring explicit role descriptions for each oversight tier \parencites[5]{6794105}. Hybrid public-sector governance case studies show \parencites[2]{0375818} that institutionalised oversight outperforms pure self-regulation: combined internal ethics boards and external regulators produced measurable reductions in bias complaints and faster appeal resolution.

Table~\ref{tab:oversight_patterns} translates these patterns into audit criteria for social recommenders. \begin{table}[H] \centering \caption{Illustrative oversight patterns and audit questions} \label{tab:oversight_patterns} \begin{tabular}{p{3.5cm}p{4.5cm}p{4.5cm}} \hline Pattern & Key features & Audit question \\ \hline HITL grading \parencites[16]{6805789} & Human validation layer, dual logging, SHAP explanations & Are recommendation overrides and rationales logged and attributable to responsible staff? \\ Hybrid DSS governance \parencites[2]{0375818} & Internal ethics board, external regulator, appeal portal & Does the platform support user appeals and independent audits of algorithmic outcomes? \\ AI Law ethical responsibility \parencites[61--62]{2977039} & Ban on manipulative and addictive systems, ethics impact reports & Are engagement optimising features assessed for digital addiction and reported in ethics disclosures? \\ HITL roles in high risk AI \parencites[5]{6794105}[7]{7461574} & Defined human in/on/over the loop responsibilities & Are human oversight roles for recommender decisions explicitly documented and resourced? \\ \hline \end{tabular} \end{table}

SHAP-based attributions in HITL grading \parencites[16]{6805789}[21]{6805789} enabled targeted overrides where model interpretation failed; the authors caution that AI alone cannot evaluate fairness or respond compassionately. Hoax-detection research reaches a parallel conclusion: low explainability undermines the legal usability of outputs \parencites[3]{9366495}. Regulators still need independent AI audits and external oversight committees to handle black-box opacity \parencites[9]{7461574}. This thesis accordingly assesses social recommender oversight through observable workflows, logging, and governance structures documented in secondary sources.

\subsection{Justice, Non-Discrimination and Inclusion}

\subsubsection{Bias, Profiling and Indirect Discrimination Risks}
\label{sec:bias_profiling_indirect_discrimination} Adaptive-education case evidence illustrates the structural pattern. Minority and non-English-speaking students are disproportionately flagged for remedial instruction because their response patterns are underrepresented in training data \parencites[5]{1308731}; rural low-connectivity students face a compounding infrastructural exclusion axis. An AI Now Institute analysis found $65\%$ of examined educational AI studies showed bias against marginalised groups \parencites[7]{1308731}, a predictable output of training pipelines that under-sample minority populations.

Sebastian's sociotechnical critique of credit-scoring warns that single-metric optimisation using demographic parity, equalised odds, or counterfactual fairness can conceal intersectional discrimination even in apparently balanced models \parencites[14]{5442231}, directing recommender audits to prioritise intersectional exposure checks; aggregate parity alone is insufficient to surface this harm \parencites[14]{5442231}[7]{1308731}. Kostenko's fairness principle gives this teeth, prohibiting models built on non-representative data or lacking bias-correction mechanisms and mandating vulnerable-group participation in development \parencites[73]{2977039}[41]{2977039}, with obligations distributed across the lifecycle through an Authorized Body for Ethics of AI maintaining a public discrimination-complaint register \parencites[41]{2977039}[42]{2977039}.

Research on ethical e-commerce shows engagement-oriented algorithms deprioritising sustainability attributes \parencites[4]{4245543}[10]{4245543}, with unverifiable green claims widening the intention-behaviour gap. The EU AI Act's unacceptable-risk tier covers psychological-vulnerability exploitation; proposed algorithmic impact assessments and demographic-performance transparency reports sharpen the regulatory stakes further \parencites[12]{6502223}[14]{5442231}[15]{5442231}. ISO/IEC bias-mitigation reference standards extend these expectations to technical design \parencites[3]{8589961}. Sebastian's finding \parencites[14]{5442231} that single-metric fixes can conceal intersectional harm underpins this report's classification of bias, profiling, and indirect discrimination as risk categories requiring documented metrics and external oversight \parencites[73]{2977039}[41]{2977039}.

\subsubsection{Protection against Social Rating and Total Scoring}
\label{sec:protection_social_rating_total_scoring} Social recommenders can drift toward quasi-universal value assessments when engagement or trust metrics are aggregated and reused across contexts. Kostenko's non-simulacrity principle explicitly prohibits this: digital twins and shadow profiles must not replace a person's autonomy or self-determination \parencites[128]{2977039}. The resulting ban covers "total scoring" that influences access to labour, credit, insurance, or justice through multi-domain data aggregation, cross-platform ID stitching, function creep, or proxy-variable discrimination.

Only narrowly targeted scoring remains permissible under five cumulative conditions: a specific lawful purpose confined to one domain, non-transferability, exclusion of protected-feature proxies, auditable transparency, and mandatory human-adjustment review \parencites[130]{2977039}. Sebastian's domain-specific fairness constraints in financial AI echo the same logic \parencites[3]{5442231}: engagement or trust scores generated within a platform cannot legitimately migrate into gatekeeping for unrelated rights or services.

Meaningful human oversight is the structural antidote. Kostenko's AI Law Model bans configurations that neutralise supervision, requiring mandatory human participation and external audits for justice, health, labour and finance decisions, with automatic denial of basic services explicitly disallowed \parencites[127]{2977039}. Event logs and appeal-rate monitoring detect "rubberstamping" \parencites[129]{2977039}. Cross-border provisions extend extraterritorial jurisdiction where AI decisions significantly affect a State's rights, security, or economy, triggering joint liability for providers, integrators, and data brokers under "flow down" due-diligence clauses \parencites[250]{2977039}. This thesis defines Code Category C: "Prohibited Total Scoring" as covering any practice that aggregates multi-domain data into a universal trust index \parencites[128]{2977039}, extends engagement scores via ID stitching into credit or hiring decisions, or maintains unreviewed blacklists derived from recommender traces \parencites[130]{2977039}. Systems operating within domain-specific, proxy-excluding, human-reviewed scoring qualify as "Targeted Scoring under Oversight" instead, though still requiring documented fairness audits \parencites[6]{7461574}[6]{4686679}[3]{5442231}. Kostenko's five-condition permissibility test \parencites[130]{2977039} is the operational dividing line between the two categories.

\subsubsection{Cultural and Linguistic Diversity in Content Curation}
\label{sec:cultural_linguistic_diversity_curation} When algorithms consistently surface dominant-language content at the expense of minority voices, the effect is structural exclusion. Kostenko's dedicated cultural-diversity principle prohibits algorithms that impose cultural stereotypes or contradict national, ethnic, or linguistic identity, requiring cultural-sensitivity mechanisms, bias testing, and complaint channels. Language support is structural non-discrimination: AI systems must support the state language plus, "within technical feasibility," minority languages without degrading translation quality or response speed \parencites[182]{2977039}.

Kostenko's broader fairness provisions reinforce this: systems trained on systematically biased data require multidimensional fairness testing and participation of diverse groups under a "fairness by design" approach \parencites[73]{2977039}[75]{2977039}.

An AI-supported mental health audit in Nigeria illustrates the tested stakes. Seven pilots used training corpora containing less than $3\%$ Fulfulde, Tiv, and Eggon content; Fulfulde proverbs were misflagged as "hostile," and Hausa female weeping was misclassified as "anger," causing PTSD under-detection measurably worse for rural low-literacy groups than urban literate users \parencites[10]{5419804}. Community Health Extension Workers acting as "algorithmic interpreters" partially mitigated the gap, translating outputs into culturally meaningful narratives through mixed audit-circle governance \parencites[11]{5419804}[2]{5419804}. AI-enabled career-counselling research reaches a parallel conclusion: responsible implementation requires localisation and cultural adaptation where large language models are trained primarily on English-language corpora \parencites{9762591}. Kostenko's algorithmic-education principle adds that AI-literacy materials must be available in multiple languages adapted to age and social group \parencites[98]{2977039}.

Gerr's diversity-index methodology using Shannon entropy and Simpson's index on Facebook, YouTube, and TikTok \parencites[3]{5142478}[5]{5142478} provides a measurement template extendable to cultural and linguistic exposure fairness. In an 800-essay HITL grading pilot, $13\%$ of AI scores were fully overridden because the model misread culturally embedded language from second-language writers \parencites[8]{5142478}[19]{6805789}; SHAP attributions identified the problematic features, and educators determined which reflected genuine weakness versus cultural encoding the model lacked training to handle.

\subsection{Wellbeing, Safety and Prevention of Harm}

\subsubsection{Psychological and Social Harm from Addictive Design}
\label{sec:psych_social_harm_addictive_design} Among $1.074$ billion Chinese short-video users, college students show diminished academic concentration, weakened social skills, and heightened anxiety and depression under excessive use \parencites[2]{8958037}. Hu links the mechanism to "dopamine driven feedback loops" produced by neural-network recommendation, micro-duration clips, infinite scrolling, and push notifications, with $22\%$ higher compulsive-checking rates in algorithmically curated feeds \parencites[2]{3303310}. Prolonged personalised-feed use further correlates with reduced grey-matter density and prefrontal activity, structurally comparable to substance dependence, and a $23\%$ higher incidence of depressive symptoms than chronological-feed controls \parencites[20]{3303310}[3]{8958037}[4]{8958037}.

Harm extends beyond the individual. Engagement-driven curation \parencites[19]{3303310} amplifies misinformation and polarisation while users remain unaware of the underlying mechanism. Gerr's European platform analysis shows personalisation systematically diminishing heterogeneous viewpoints \parencites[29]{5142478}. Instagram's Eurocentric beauty-standard prioritisation is linked to increased body dissatisfaction among adolescent women of colour \parencites[3]{3303310}. The DSA's TikTok investigation singles out infinite scroll and autoplay as potential manipulative design \parencites[2]{8958037}, and the AI Act's unacceptable-risk tier \parencites[2248]{6502223} bans systems exploiting psychological vulnerabilities; the anticipated Digital Fairness Act \parencites[390]{7896734} extends this to addictive product design, moving the constraint toward enforceable obligation.

Kostenko's ethical-responsibility principle bans systems causing digital addiction or suppressing autonomy \parencites[61--62]{2977039}, with confidentiality provisions barring profiling-driven advertising that restricts free choice without explicit consent \parencites[45]{2977039}. Dark-pattern research on Indian e-commerce confirms that default bias, social pressure, and friction-loaded exits operate through interface structure alone \parencites[8]{6295677}[7]{6295677}; the Digital Fairness Act \parencites[390]{7896734} is expected to codify this as a first-class enforcement category.

Hu's Algorithmic Impact Assessment Scale (AIAS-5) quantifies emotional volatility, sleep disruption, and compulsive usage, feeding a proposed "digital nutrition label" disclosure regime and dynamic transparency index \parencites[20]{3303310}[19]{3303310}. Corporate-governance research adds independent audit trails as a mechanism for surfacing ethical liabilities \parencites[12]{8045684}. Engagement strategies materially elevating AIAS-style dimensions, especially combined with opaque profiling and no meaningful exit option, fall into high-risk or prohibited zones in this thesis \parencites[2248]{6502223}[390]{7896734}.

\subsubsection{Risk Thresholds and Unacceptable Practices}
\label{sec:risk_thresholds_unacceptable_practices} The EU AI Act's "unacceptable risk" tier explicitly names systems exploiting psychological vulnerabilities alongside social scoring \parencites[2248]{6502223}. Kostenko's AI Law Model draws the same boundary, prohibiting behaviour manipulation and digital addiction as a mandatory certification criterion \parencites[61--62]{2977039}. Together these define Code Category C, "Unacceptable Manipulative Practices": designs intentionally exploiting psychological vulnerabilities \parencites[2248]{6502223}[61--62]{2977039}, addictive infinite-scroll and variable-reward mechanics without meaningful exit, and architectures blocking the DSA-mandated non-profiled feed option. The anticipated Digital Fairness Act's targeting of addictive design \parencites[390]{7896734} adds a further enforcement layer, extending platform liability beyond the AI Act and DSA simultaneously.

A second ceiling covers discrimination and social rating. Sebastian's pre-deployment fairness verification research warns that single-metric fixes mask intersectional harm \parencites[14]{5442231}. Kostenko's fairness principle bans social-scoring "integral assessments" affecting service access \parencites[41]{2977039}. Cross-domain trust indices derived from engagement traces or unmitigated disparate content exposure fall under Code Category C: "Unacceptable Scoring and Discrimination" \parencites[41]{2977039}[14]{5442231}.

Operational thresholds are achievable between prohibited and compliant practice. The Algorithmic Impact Navigator \parencites{2742041} documents one case: a $30\%$ lending-processing-time reduction remained fair where denial gaps were statistically insignificant under routine audits and continuous monitoring. Sebastian's domain-specific fairness constraints \parencites[14]{5442231} supply the design logic: hybrid human-AI governance with compulsory authorisation for sensitive decisions built into workflows is the configuration this thesis classifies as Code Category B, "High Risk but Potentially Mitigable."

Unacceptable practice is as much a governance failure as a design one. A Nigerian mental-health AI trial improved PTSD remission under Pastoral Data Trust governance \parencites{5419804}, yet also documented a false-positive suicide alert triggering a violent vigilante patrol and a privacy breach from bypassed controls: nominal metric compliance did not prevent documented harm. Conservative thresholds and enforceable misuse sanctions are therefore non-negotiable. Findings combining high engagement optimisation with weak governance are classified Code Category C in this thesis regardless of nominal compliance scores \parencites[61--62]{2977039}. ISO/IEC risk-management standards and AI Act/GDPR compliance research \parencites{8589961}[1753]{7461574}[2248]{6502223}[390]{7896734} extend these expectations globally.

\subsubsection{Duty of Care for Vulnerable and Minor Users}
\label{sec:duty_of_care_vulnerable_minors} Kostenko's digital-safety-of-the-child principle requires systems processing children's data to be "secure by default" across the lifecycle, with developers maintaining Child Impact Assessments and Child Safety Data Sheets, and regulators entitled to full technical dossiers \parencites[214]{2977039}. Explainability standards are explicitly age-sensitive: plain-language decision-logic interpretations and children's-rights tools such as profiling refusal are mandatory, and manipulative or obscure interfaces are classified as violating the child's best interests, with suspension authority available as a remedy \parencites[214]{2977039}.

Age-appropriate design provisions translate these principles into observable interface constraints. Minors' defaults must minimise data collection and disable profiling. Endless scrolling, autoplay, and forced-exit penalties are prohibited. "No personalisation" switches and "why am I seeing this" explanations are mandatory \parencites[223]{2977039}. Hu's $23\%$ higher incidence of depressive symptoms in algorithmically curated versus chronological-feed controls \parencites[20]{3303310} is a particular measure of exposure risk for a group that lacks fully developed inhibitory control, making these interface constraints directly relevant to the short-video mechanics documented above.

Protection extends to profiling and monetisation. Psychotype- or emotion-based targeted advertising to minors is banned \parencites[20]{2977039}[224]{2977039}; predictive "life path" profiling is outlawed, and data resale or third-party model training requires a separate best-interests legal basis \parencites[223]{2977039}. Business models premised on behavioural advertising to minors conflict with this baseline, and the DSA's systemic-risk framework along with the UK Online Safety Act both impose additional child-content obligations \parencites[2248]{6502223}[2249]{6502223}. UNESCO's AI Ethics Recommendation calls for algorithmic diversity to ensure pluralism and positions platforms as custodians of psychological welfare, extending these obligations beyond European law \parencites[2249]{6502223}.

Governance architecture beyond UX adjustments is required. Kostenko mandates independent ethics commissions with suspension authority and annual AI Ethics Impact Reports \parencites[61--62]{2977039}. Corporate human-AI audit frameworks show clearly defined roles improve detection power without sacrificing professional judgment \parencites[2]{6539584}[3]{6539584}. Platforms combining intensive profiling, addictive mechanics, and undocumented child-safety dossiers fall into "Code Category C: Prohibited or Critical Risk"; secure-default, age-appropriate, human-audited configurations exemplify duty-of-care-compliant design \parencites[214]{2977039}[223]{2977039}[16]{8200246}. Regulators should also scrutinise how facial and interactional data are repurposed for psychometric targeting of children, since innocuous signals can be recombined to reveal sensitive traits with documented accountability chains required at every data-handoff stage \parencites[4]{5750298}.

\section{Qualitative Methodology and Audit Framework}

\subsection{Overall Qualitative Research Design}

\subsubsection{Choice of Secondary Qualitative Synthesis}
\label{sec:choice_secondary_synthesis} Proprietary recommender logs are unavailable; that structural constraint drives the choice of secondary qualitative synthesis. Rich documentary material exists across regulatory, managerial and ethical domains. Dark-pattern research on Indian e-commerce justifies this design where "primary experimentation or access to proprietary platform data is limited" \parencites[6]{6295677}. Asante and Adeoba's Reflexive Thematic Analysis spanning marketing, information systems and CSR literature \parencites[8]{8200246}[10]{8200246} confirms that the same analytic backbone transfers across domains, as does Malik's thematic synthesis of policy and practice sources for AI-enabled career counselling \parencites{9762591}. Systematic coding generates analytical themes around addiction, manipulation and auditability, with every coding decision explicitly documented. Line-by-line coding and multi-coder CASP appraisal used in clinical-ethics thematic synthesis \parencites[7]{6419865} supply the procedural template.

Understanding how the AI Law Model and DSA-style regimes articulate risk-management duties is a conceptual goal, addressed by doctrinal-empirical means. The hybrid approach Khare et al. \parencites[6]{6295677} use to evaluate regulatory adequacy under proprietary secrecy is the closer methodological analogue. Quality safeguards follow Lincoln and Guba's trustworthiness criteria as applied by Asante and Adeoba: methodological description, audit trails and confirmability \parencites[10]{8200246}. This design avoids direct human-subjects engagement, relying on regulatory reports, judicial decisions and complaint statistics. Kostenko's AI Law Model \parencites[54]{2977039} calls for Human Rights Impact Assessments and incident logs grounded in ISO/IEC 23894 and the NIST AI RMF, making thematic synthesis of textual artefacts both a methodological necessity and thematically consistent with the audit-toolkit target \parencites[12]{2977039}[23]{6839355}.

\subsubsection{Alignment between Research Questions and Methods}
\label{sec:alignment_rq_methods} Three research questions, fixed since the approved Capstone exposé (see Appendix A), drive this project. Each addresses conceptual relationships rather than causal magnitudes: how engagement-optimised recommenders instantiate manipulative design and algorithmic addiction; how EU-style regulation frames duties and prohibitions; and how an Ethical AI Audit Toolkit can be constructed using CSR-based organisational governance. Secondary qualitative synthesis fits this configuration because the aim is to map patterns and interpret mechanisms. Khare et al. express the same goal: "map patterns, interpret behavioural mechanisms, and evaluate regulatory adequacy" \parencites[6]{6295677}. Asante and Adeoba similarly use secondary synthesis \parencites[8]{8200246}[10]{8200246} to theorise how Communities of Practice and CSR mediate ethical AI implementation.

Gerr's convergent mixed-method design pairing DSA/AI Act policy-document analysis with bias and diversity metrics on Facebook, YouTube and TikTok \parencites[3]{5142478} provides a methodological anchor for the policy-document coding used here. Coding transparency draws further on the MMR-RHS checklist's emphasis on documenting sampling, integration logic and credibility strategies \parencites[2]{8591813}, and on the grounded-theory triangulation practices Kishore et al. \parencites[6]{4245543} use for AI-mediated finance. Absent new interviews, the same coding logic is applied to documentary material, with that substitution declared as an explicit limitation. Research questions are formulated at the level of observable behaviour and documented harm, excluding reverse engineering or telemetry collection. Asante and Adeoba and Khare et al. both work with published accounts alone \parencites[11]{8200246}[6]{6295677}. Mixed-methods mental-health research using an ethnographically co-designed "CHEW Competency Wheel" \parencites[4]{5419804} further supports framing questions around AI interaction with vulnerable communities through qualitative design \parencites[6]{6295677}[8]{8200246}.

\subsection{Data Sources and Selection Strategy}

\subsubsection{Regulatory and Policy Texts as Data}
\label{sec:reg_policy_texts_data} Regulatory and policy documents are primary empirical data, not background. Khare et al. code ASCI guidelines, CCPA advisories and enforcement orders such as BookMyShow \parencites[6]{6295677} directly as evidence of normative expectations and governance performance. Gerr treats DSA, the German Action Plan and EU AI Act texts the same way, coding them alongside platform policies for transparency and pluralism themes \parencites[3]{5142478}. Kostenko's AI Law Model is analytically rich data by the same standard. Its technological-sovereignty principle allows states to prohibit AI systems where independent audit is impossible or critical dependence exists \parencites[107]{2977039}, and its generative-AI regime mandates registration, labelling and independent audits covering manipulative impact and bias \parencites[31]{2977039}.

DSA analyses supply further discrete taxonomy variables. Very large platforms must conduct systemic risk assessments, offer a non-profiled feed option, and submit to annual audits, with Article 40 granting researcher data access \parencites[7]{1456534}[19]{3195574}. The anticipated Digital Fairness Act targets addictive design and unfair personalisation, aiming to close gaps between the DSA, DMA and unfair-commercial-practices law \parencites[390]{7896734}. "Non-profiled feed requirement" is one such variable extracted directly from these texts. The EU AI Act's four-tier risk scheme, classifying systems exploiting psychological vulnerabilities as unacceptable risk, anchors Code Category C \parencites[12]{6502223}. Kostenko's explicit bans on manipulative and addictive systems reinforce that anchor \parencites[61--62]{2977039}[107]{2977039}. Regulatory texts also encode implementation-capacity limits. Research on Ukraine's "imitative modernization" of GDPR/AI Act/DSA frameworks documents a "double imitation" gap between formal transposition and effective governance \parencites[19]{5162643}. Literacy surveys show high self-assessed digital competence but low ability to recognise manipulation, alongside scepticism about cross-border enforcement \parencites[10]{9360054}. Both are coded here as contextual constraints on any audit toolkit that assumes strong supervisory capacity.

\subsubsection{Academic and Grey Literature on Algorithmic Addiction}
\label{sec:academic_grey_addiction} Two neuropsychological strands anchor coding of algorithmic addiction. Li's synthesis on short-video platforms identifies a clear mechanism: neural-network recommendation combined with 15--60 second clips, infinite swipe navigation, autoplay and push notifications creates what Li terms an "addiction loop," with Chinese usage data averaging near $60$ minutes and $21$+ sessions daily, especially among students \parencites[2]{8958037}[3]{8958037}. Dopamine surges occur before each swipe, structurally akin to slot-machine variable rewards. Chronic stimulation elevates reward thresholds \parencites[3]{8958037}[4]{8958037}. Hu adds quantitative weight: algorithm-driven feeds produce $22\%$ higher compulsive-checking rates, link to depressive symptoms, and are associated with reduced prefrontal activity and reduced grey-matter density paralleling substance-dependence patterns \parencites[2]{3303310}[20]{3303310}. Kishore et al.'s grounded-theory study of AI-embedded financial tools identifies a second harm pathway: users treat algorithm outputs as advisory and frequently override them, flagging automation bias and diminished agency \parencites[7]{4245543}[6]{4245543}.

Methodological precedents justify synthesising these diverse sources without primary data collection. A PRISMA-based review of AI-powered mobile proctoring screened 180 records and synthesised 20 on surveillance anxiety and algorithmic bias \parencites[2]{2055956}[9]{2055956}. A structured review of AI in social innovation \parencites[3]{4246340} used inductive three-stage coding to identify bias and an "accountability void." A conceptual synthesis on AI-enabled career counselling derived governance clusters through secondary sources alone \parencites[14]{9762591}. Grey and model-law texts are treated as data on the same basis: Kostenko's AI Law Model prohibits systems causing digital addiction or suppressing autonomy, articulating a right to a psycho-emotionally safe digital environment \parencites[61--62]{2977039}[119]{2977039}, coded here as a normative threshold informing Code Category C. Indian dark-pattern studies show that manipulative interfaces exploiting dual-process cognition can be documented through consumer surveys and enforcement records alone \parencites[6]{6295677}[7]{6295677}. Coding across all these sources distinguishes direct empirical measures from inferred neurocognitive mechanisms and normative governance responses.

\subsubsection{Public UX Analyses and Advocacy Reports}
\label{sec:public_ux_advocacy_reports} Roughly $87\%$ of approximately $30$ high-impact Indian platforms deployed at least one dark pattern. Khare et al.'s audit \parencites[7]{6295677} is treated throughout as primary empirical data. Their coding against established HCI taxonomies maps interface interference, obstruction, subscription traps and privacy zuckering onto System 1 biases including default bias, scarcity and social proof, and their "sludge"-based obstruction construct directly informs working codes such as "false urgency widgets" and "roach motel subscription flows" \parencites[7]{6295677}. EU regulatory advocacy supplies a second type of UX evidence. Platform X's €120 million fine for deceptive verification features and Temu's DSA breach investigation over addictive design and false urgency cues \parencites[390]{7896734} inform "Code Category C: deceptive verification and cancellation flows." The Georgian Competition and Consumer Agency's sanctioning \parencites[390]{7896734} of false countdown timers, fake reviews and deceptive brand imagery (the "Maxbosch"/BOSCH case) adds further documented UI artefacts to the coding matrix, confirming that these patterns are regulator-identified harms.

A subset of UX literature addresses AI-mediated environments. Mogoi et al.'s review of AI proctoring \parencites{2055956} records user perceptions of being "constantly watched," linking facial recognition and gaze tracking to surveillance anxiety, with privacy-by-design and transparent communication treated as UX governance benchmarks. HITL grading research shows that explicit AI/human labelling and accessible appeal procedures increase student trust \parencites{6805789}, coded here as positive UX patterns contrasting with e-commerce dark patterns. These sources anchor abstract concepts like "manipulative interface" in observable structures, following Khare et al.'s content analysis of screenshots and flows \parencites[7]{6295677}[10]{8200246}. Internal algorithm logic is not directly observable; this analysis works through documented UI behaviour and enforcement decisions alone.

\subsection{Development of a Working Coding Matrix}

\subsubsection{Code A Interface Mechanics and Engagement Features}
\label{sec:codeA_interface_mechanics} Code A covers observable interface mechanics implicated in algorithmic addiction. Every slot is coded at the level of explicit UI elements; proprietary algorithms are assessed via published behaviour descriptions \parencites[6]{6295677}[2]{8958037}, not code or telemetry. Three clusters structure the scheme.

The first cluster addresses high-frequency low-friction engagement loops. Li's "addiction loop" of infinite swipe feeds, autoplay and push notifications shifts users into "auto pilot" scrolling averaging near $60$ minutes and $21$+ sessions daily \parencites[2]{8958037}[3]{8958037}. Hu links those same mechanics to $22\%$ higher compulsive checking and reduced prefrontal activity \parencites[2]{3303310}[20]{3303310}. These findings motivate subcodes A1 (infinite scroll), A2 (autoplay) and A3 (high-frequency push notifications). The second cluster addresses salience and urgency constructs. Countdown timers, "only two left" stock alerts and visually dominant acceptance prompts \parencites[136]{8745065}[7]{6295677} motivate A4 (countdown/scarcity banners) and A5 (visual hierarchy favouring acceptance over decline). The third cluster captures transparency-supporting mechanics: HITL grading's explicit AI/instructor labelling and override visibility \parencites[20]{6805789}[21]{6805789}, and the AI Law Model's requirement \parencites[59.1]{2977039} that AI outputs be labelled and logged with model identifiers, together motivate A6 (AI labelling), A7 (non-personalised mode toggles) and A8 (human-validation indicators). These codes originate in educational contexts but apply to social recommenders because their presence signals whether engagement optimisation is coupled to meaningful oversight. Research on AI-driven review systems shows that explicit labels and audit logs help restore user trust \parencites[4]{7235975}. Transparency affordances are coded separately from inferred algorithmic behaviour because their presence or absence is directly observable from published interface descriptions.

Table~\ref{tab:codeA_slots} summarises Code A slots. \begin{table}[H] \centering \caption{Code A slot definitions for interface mechanics and engagement features} \label{tab:codeA_slots} \begin{tabular}{p{1.8cm}p{5.0cm}p{5.0cm}} \hline Slot & Description & Key sources \\ \hline A1 & Infinite scroll or continuous feed without natural stop cues & Short video "addiction loop," persistent scrolling \parencites[2]{8958037}[3]{8958037}; dopamine driven checking \parencites[2]{3303310} \\ A2 & Auto play of next item without user click/tap & Auto play in short video architectures \parencites[2]{8958037} \\ A3 & High frequency, behaviour based push notifications & Push notifications within addiction loop \parencites[2]{8958037}[3]{8958037} \\ A4 & Countdown timers or scarcity banners in decision flows & Urgency marketing cues \parencites[136]{8745065}; false urgency patterns \\ A5 & Visual or interaction hierarchy favouring acceptance over refusal & Interface interference and obstruction \parencites[7]{6295677} \\ A6 & Explicit labelling of AI generated outputs & Labelling requirements in education \parencites[59.1]{2977039}; dual sourced feedback UI \parencites[20]{6805789} \\ A7 & User accessible non personalised or "off" modes & Non profiling recommendation option under DSA \parencites[7]{1456534}; age appropriate "no personalisation" switches \parencites[223]{2977039} \\ A8 & Visible human validation or override indicators & HITL grading override logs and authority retention \parencites[20]{6805789}[59.1]{2977039} \\ \hline \end{tabular} \end{table}

Coding follows open-then-closed slot assignment for inter-coder consistency, similar to YAML-based codebooks used in large language model prompting research \parencites{7286763}{8584408}, validated by cross-checking Li's and Hu's descriptions against DSA enforcement concern over manipulative design impairing free decision-making \parencites[2]{8958037}[2248]{6502223}. High-level fairness and transparency principles are a useful normative frame but require actionable operational rules to apply reliably across platform designs \parencites[3]{2168844}. One caveat applies broadly: models jointly considering temporal dynamics and social trust \parencites[2]{0335538} suggest that what appears as an addictive A1--A5 pattern in one cohort could reflect shifting preference dynamics in another, so auditors should record cohort characteristics alongside code assignments.

\subsubsection{Code B Psychological Triggers and Motivational Effects}
\label{sec:codeB_psych_triggers} Code B records psychological and motivational mechanisms exploited by interface patterns. Every slot is coded from triggers explicitly documented in empirical and regulatory sources; inferred design intent alone is not sufficient grounds for coding.

Time pressure and scarcity form the first cluster. Countdown timers, "only two left" notices and expiring discounts reframe routine purchases as time-sensitive decisions targeting budget-constrained consumers \parencites[136]{8745065}. Indian dark-pattern audits formalise this as "false urgency/scarcity" \parencites[7]{6295677}, yielding B1 (time-pressure cues) and B2 (scarcity/loss framing). Cognitive load and decision fatigue form the second cluster. "Choice obstruction" describes subscription traps and roach motels that keep users in fast, heuristic-driven decision modes \parencites[6]{6295677}. Urgency-marketing interviews separately document "decision fatigue" and "psychological targeting" \parencites{8745065}, yielding B3 (friction-based obstruction) and B4 (decision-fatigue amplification). The third cluster covers temporal orientation. Short-video platforms' "just one more swipe" aligns with present-oriented focus and undermines long-term goals \parencites[2]{8958037}[3]{8958037}. Hu's $22\%$ compulsive-checking and $23\%$ depressive-symptom findings capture the shift toward immediate algorithmic reward \parencites[2]{3303310}[20]{3303310}, yielding B5 (instant reward salience) and B6 (temporal discounting bias). A fourth cluster addresses social and identity triggers. Fabricated social proof, "confirm shaming" \parencites[7]{6295677} and peer/influencer-driven guilt documented in grounded-theory consumption research \parencites{4245543} yield B7 (social proof/conformity pressure) and B8 (guilt/shame messaging).

Gamified-system research provides a positive-coding counterpoint. Intrinsic-motivation taxonomies applied to avatars, rewards and pacing show genuine autonomy/competence support when aligned with educational aims \parencites{5645864}. Token-economy research shows that "technostress" from complexity can cancel those gains \parencites{9283805}, yielding B9 (autonomy-supportive feedback loops) and B10 (technostress-induced demotivation), distinguishing designs that reinforce motivation from those that overload users. A final boundary condition comes from affective computing: emotional-AI systems risk "behavioral interaction addiction" and "social relationship alienation" absent transparency and assistive, non-anthropomorphic design \parencites{9738431}, yielding B11 (emotional dependence triggers), coded where systems simulate companionship without signalling artificiality. Table~\ref{tab:codeB_summary} summarises Code B slots and anchor sources. \begin{table}[H] \centering \caption{Code B psychological triggers and motivational effects} \label{tab:codeB_summary} \begin{tabular}{p{1.1cm}p{5.0cm}p{5.0cm}} \hline Slot & Description & Key sources \\ \hline B1 & Time pressure cues (countdown timers, expiring discounts) & Urgency marketing cues \parencites[136]{8745065} \\ B2 & Scarcity and loss framing messages (stock notices exaggerating availability) & False urgency / scarcity patterns \parencites[7]{6295677} \\ B3 & Friction based obstruction (subscription traps, roach motels, complex cancellations) & Choice obstruction \parencites[6]{6295677} \\ B4 & Decision fatigue amplification through design induced cognitive exhaustion & Decision fatigue and psychological targeting \parencites{8745065} \\ B5 & Instant reward salience from swipe or scroll interactions & Short video addiction loop \parencites[2]{8958037}; compulsive checking \parencites[2]{3303310} \\ B6 & Temporal discounting bias favouring present oriented, low cost interactions & Present oriented focus undermining long term goals \parencites[3]{8958037}[20]{3303310} \\ B7 & Social proof and conformity pressure (fabricated or amplified testimonials) & Confirm shaming and generative social proof \parencites[7]{6295677} \\ B8 & Guilt or shame based messaging exploiting emotional or identity cues & Peer influence and post purchase regret \parencites{4245543} \\ B9 & Autonomy supportive feedback loops in gamified systems & Intrinsic motivation via avatars and pacing \parencites{5645864} \\ B10 & Technostress induced demotivation from complexity or instability & Token economies and technostress \parencites{9283805} \\ B11 & Emotional dependence triggers from simulated companionship or reciprocity & Behavioural interaction addiction in affective computing \parencites{9738431} \\ \hline \end{tabular} \end{table}

\subsubsection{Code C Legal Risk Levels and Compliance Status}
\label{sec:codeC_legal_risk_levels} Code C assigns legal risk by reference to explicit prohibitions and documented duties. Three binding or model-law sources provide the anchors: the AI Law Model's bans on manipulative and non-auditable systems \parencites[68.5]{2977039}[70]{2977039}, its TSAI provisions on social rating and covert censorship, and the EU AI Act's unacceptable-risk classification for systems exploiting psychological vulnerabilities \parencites[31.4]{2977039}.

The AI Law Model prohibits systems that cannot be independently audited, lack national control over critical components, or embed hidden behavioural-influence mechanisms, requiring placement in a register of prohibited AI \parencites[31.4]{2977039}. TSAI provisions separately ban social rating, manipulative political targeting and covert censorship \parencites[68.5]{2977039}. Neurorights research documents how covert influence operates through interface and reward-schedule mechanics that bypass conscious deliberation \parencites[11]{0783007}. Configurations matching these patterns are assigned to "Code Category C: Prohibited Systemic Practices," signalling risk not mitigable by incremental adjustment. A second rule cluster defines non-negotiable traceability conditions. High-risk systems must maintain decision-architecture documentation, full audit trails linking inputs to outputs, and training-data bias metrics as part of a technical passport accessible to state bodies \parencites[19.5]{2977039}. TSAI further requires unique identifiers, algorithmic passports and automated suspension capability for cross-border systems \parencites[68.6]{2977039}. Absence of these artefacts in high-impact recommendation contexts is coded "Critical Traceability Failure." Digital-sovereignty provisions add a territorial dimension: localisation requirements prevent "digital colonization" by foreign-controlled critical infrastructure, envisaging restrictions where public registers or decisions rely on government-inaccessible code \parencites[33.2]{2977039}[33.4]{2977039}. Recommenders routing governance functions through opaque foreign infrastructure fall into "high or unacceptable" risk bands accordingly.

Table~\ref{tab:codeC_levels} summarises these working legal risk levels. \begin{table}[H] \centering \caption{Code C working legal risk levels and compliance cues} \label{tab:codeC_levels} \begin{tabular}{p{3cm}p{4.5cm}p{4.5cm}} \hline Code C category & Indicative practices & Key legal anchors \\ \hline Prohibited systemic practices & Social rating indices, manipulative political targeting, covert censorship, non auditable AI & TSAI prohibitions on social rating and manipulative targeting; auditability bans \parencites[31.4]{2977039}[68.5]{2977039} \\ Critical traceability failure & No technical passport, missing logs, no event evidence for decisions in high risk settings & Documentation and logging duties for high risk systems \parencites[19.5]{2977039}[68.6]{2977039} \\ High sovereignty risk & Critical decisions based on foreign controlled servers, inaccessible code or data relocation & Digital sovereignty, local storage and anti colonization clauses \parencites[33.2]{2977039}[33.4]{2977039} \\ Statutory compliance with oversight & Registered systems, annual audits, public trust labelling, accessible registry APIs & TSAI supervision, auditing and trust marks \parencites[68.7]{2977039} \\ \hline \end{tabular} \end{table}

Code C can also mark positive compliance. Healthcare and civil-service HR systems under statutory or co-regulatory models show continuous bias audits and appeal mechanisms, producing measurable bias reduction and faster appeal resolution \parencites[4]{0375818}, mapped here to "Statutory Compliant with Effective Oversight." Revenue-process audit research reinforces that automation architectures require explicit automated/human-component documentation and feedback loops for anomaly detection \parencites{6539584}. Practical recommender-audit guides stress that platform-profile and stakeholder mapping are essential first steps. Absent such documentation, high-impact social recommenders shift from "moderate" into higher risk categories \parencites[70]{2977039}[4]{0375818}{6539584}.

\subsection{Methodological Limitations and Quality Criteria}

\subsubsection{Constraints of Observing Black-Box Algorithms Indirectly}
\label{sec:constraints_black_box_observation} Only outward behaviour is observable. That constraint is structural, matching Khare et al.'s reliance "exclusively on secondary data" where "primary experimentation or access to proprietary platform data is limited" \parencites[6]{6295677}. The contrast with internal access is instructive: embedded-audit pilots on a mental-health triage model report subgroup sensitivity metrics down to dialect-level error patterns \parencites[9]{5419804}, and HITL grading research reports override rates and criterion-level $F1$ scores because the dashboard was under researcher control \parencites[17]{6805789}[21]{6805789}. Neither type of access is available for commercial social recommenders.

Regulatory texts assume that internal observability exists. The AI Law Model demands technical passports with full input-output tracing available to authorities \parencites[71]{2977039}. ISO/IEC standards for bias assessment aim to reduce black-box opacity \parencites[2]{8589961}. External scholars can typically verify only that a standard is claimed; its implementation in a specific pipeline remains opaque. Three linked limitations follow. Systematic ranking bias cannot be quantified at model level; the analysis relies on regulators' findings and UX audits that carry their own sampling and jurisdictional constraints \parencites[6]{6295677}[3]{5142478}. Compliance with documentation duties central to the AI Law Model remains unobservable unless platforms voluntarily publish audit summaries \parencites[71]{2977039}[51]{2977039}. Causal attribution from interface mechanics to individual harm is mediated through third-party research using context-specific samples \parencites[9]{5419804}. Khare et al.'s explicit position that such designs "map patterns" and do not "establish causal relationships" \parencites[6]{6295677} defines the appropriate scope. This thesis analyses algorithmic addiction through documented patterns in external audits and regulatory texts, treating Kostenko's safety-and-reliability provisions on incident recording and independent verification \parencites[38]{2977039} and procedural-safeguard proposals for emergency shutdown and human override \parencites[86]{1225951} as normative benchmarks whose empirical status has not been independently verified.

\subsubsection{Reliability and Transparency in Qualitative Audit Work}
\label{sec:reliability_transparency_qual_audit} Systematic procedures carry the reliability burden here. Asante and Adeoba's trustworthiness table expresses credibility through triangulation and confirmability through reflexive interpretation \parencites[10]{8200246}[11]{8200246}. Kishore et al. add open/axial/selective coding with analyst triangulation and saturation checks for AI decision-support research \parencites[6]{4245543}. Debt-behaviour research contributes member checking and coding audit trails \parencites[8]{7042289}. This project's reliability baseline follows these precedents: explicit coding schemas, saturation checks, documented analytic decisions, and verbatim coding procedures with multi-researcher dispute resolution \parencites[4]{5716309}[5]{5716309}.

Transparency operates through the same mechanism Khare et al. use, describing their $\sim$30-platform sampling frame and HCI-typology coding explicitly \parencites[7]{6295677}. The working coding matrix (Codes A, B, C) is explicitly defined and never presented as official regulatory classification. An Algorithmic Governance Assessment Matrix for zakat institutions \parencites[59]{2702055} shows how secondary quantitative indicators integrate into qualitative audit while keeping observed values distinct from interpretive conclusions. That distinction is maintained throughout between reported metrics and categorisations like "Code Category C." Explainable-AI research on hoax detection argues that black-box classification without explanation undermines public trust and legal usability; audit findings must trace to specific UI features or documented metrics \parencites[2]{9366495}. Multimodal emotion-recognition research sets a realistic bound: high classification performance can introduce steering vectors difficult to trace from public reports, limiting what can be inferred from platform disclosures \parencites[3]{3625351}. The AI Law Model embeds transparency structurally, requiring technical documentation and cross-border liability rules guaranteeing supervisory access \parencites[107]{2977039}[111]{2977039}, and published platform reports are assessed against those criteria even without code inspection.

\subsubsection{Reflexivity and Researcher Positionality}
\label{sec:reflexivity_positionality} Secondary qualitative synthesis does not produce findings from neutral ground. Position shapes interpretation. Asante and Adeoba ground exactly this in constructivism: conceptual understanding emerges through interpretation of existing evidence, not direct participant engagement \parencites[8]{8200246}[10]{8200246}. The analysis here is unavoidably theory-laden, filtered through prior commitments about AI governance and CSR. Selective reading is a real risk, mitigated by the analyst-triangulation and audit-trail practices used in constructivist grounded theory and reflexive thematic analysis \parencites[6]{4245543}[10]{8200246}[11]{8200246}.

Interpretive bias is especially sensitive where AI intersects with ethical or religious values. Research on Islamic AI applications \parencites[5]{0140790} shows how researchers explicitly confronted their own positionality, maintained reflexive logs and consulted scholars from diverse orientations to resolve coder disagreements. A comparable risk exists here given this author's professional orientation toward AI ethics and a tendency to treat precautionary regulatory proposals as inherently desirable. Reflexive memoing and inclusion of divergent capacity assessments \parencites[19]{5162643}[10]{9360054} guard against that tendency. This thesis privileges regulatory documents and enforcement records for interface-manipulation questions, following Khare et al. \parencites[6]{6295677}, and peer-reviewed or policy literature for governance questions, following Asante and Adeoba \parencites[8]{8200246}. That choice embeds a judgement that formal law is authoritative for defining acceptable recommender behaviour; it may under-represent governance imaginaries outside the EU model-law tradition \parencites[12]{6502223}[22]{2977039}.

How black-box opacity is read also reflects this orientation. Mixed-methods AI-mental-health research stresses human-mediated interpretation to avoid "AI mental health colonialism" \parencites[11]{5419804}, and AI-proctoring research treats institutional readiness as equally critical to algorithmic accuracy \parencites[9]{2055956}. Both reinforce this thesis's focus on governance architecture. Kostenko's AI Law Model, treated throughout as a normative benchmark prioritising human rights and independent oversight \parencites[22]{2977039}[37]{2977039}[44]{2977039}, and Asante and Adeoba's framing of CSR as the natural locus of reform \parencites[11]{8200246}, together shape the coding matrix's assumptions. An innovation-first standpoint would weight the same evidence toward different risks; readers should bear that interpretive latitude in mind \parencites[2]{0375818}[7]{7461574}.

\section{Comparative Platform Audit}

\subsection{Selection Rationale and Platform Profiles}

\subsubsection{Short-Form Video Platforms with Personalized Feeds}
\label{sec:shortform_personalized_feeds} With $1.074$ billion short-video users in China alone ($95.4\%$ of netizens), Li's description of TikTok-type feeds as an "addiction loop" \parencites[2]{8958037} captures the architecture directly: infinite scroll removes stopping points, autoplay removes the need to choose, and push notifications recapture users who left. Hu traces the measurable consequences: "dopamine driven feedback loops," $22\%$ higher compulsive checking, reduced prefrontal activity, and $23\%$ higher depressive-symptom incidence \parencites[2]{3303310}[20]{3303310}. The DSA's 2026 TikTok proceeding names these features as potential "systematic manipulative design" \parencites[2]{8958037}[15]{6502223}, and Chitra et al. document the same recommender driving viral misinformation with community guidelines as the only visible restraint. Short-video addiction correlates with impaired sleep and "information cocoons" that attenuate critical thinking \parencites[6]{8958037}[8]{6502223}, while politically interested accounts are "quickly directed toward extremist material." Kostenko's digital-ecology provisions ban algorithms that "form addiction or create self-replicating digital environments" eroding autonomous thought \parencites[119]{2977039}. Short-form personalised video sits at the highest-priority intersection of Code A, Code B and Code C in this report's working taxonomy \parencites[2]{8958037}[20]{3303310}[119]{2977039}[15]{6502223}[2]{3303310}[6]{6295677}.

\subsubsection{Image and Story-Centric Social Networks}
\label{sec:image_story_social_networks} Instagram's "Explore" page amplified body-negative imagery by $23\%$ relative to chronological feeds; algorithmically prioritised content shows $31\%$ higher negative affect than baseline \parencites[3]{3303310}, with statistically significant self-esteem declines among adolescent women of colour. Quinelato classifies these platforms as "social media" that algorithmically reference third-party content while collecting granular interaction data for profiling \parencites[4]{6941611}; Ferraro et al. describe recommendation-driven consumption as "monitoring-based marketing" where each interaction doubles as ranking-logic input \parencites[5]{6314747}. Regulatory scrutiny focuses on behavioural advertising, treating vulnerable demographics profiled via third-party tracking as the primary risk population, with opaque persuasion as the operative concern \parencites[4]{6941611}. Literacy surveys sharpen the picture: users rate themselves confident in routine digital tasks yet score only a mean $2.31/5$ on recognising manipulation \parencites[10]{9360054}. Fairness in these systems is ongoing, requiring continuous auditing; documentation practices proven in clinical and criminal-justice AI apply equally to revealing representational gaps in image-based recommendation \parencites[20]{2616018}. The Algorithmic Impact Navigator's fairness and explainability requirements \parencites[2]{2742041} and the zakat-distribution governance matrix's transparency dimensions \parencites[4]{2702055} are operative wherever AI-intensive architectures produce documented psychosocial impacts at this scale \parencites[3]{3303310}[5]{6314747}.

\subsubsection{Professional Networking Platforms with Algorithmic Feeds}
\label{sec:prof_networking_algo_feeds} Professional networking platforms with algorithmic feeds warrant inclusion in this audit, but a firm evidentiary limit must be stated first. None of the sources gathered for this report name LinkedIn specifically. Claims in this subsubsection are extrapolated from adjacent recommender-driven platforms and are marked illustrative rather than platform-verified, in line with this thesis's Integrity Rule 1.

Professional feeds combine engagement mechanics with economic gatekeeping around jobs and reputation, stakes that exceed those of entertainment feeds. Chitra et al. list Facebook, Instagram, TikTok, X and Google Search as nodes where recommenders are "invisible gatekeepers" in political information flows \parencites[2]{6502223}; the same engagement-ranking risks transfer by extrapolation to professional contexts. Gerr shows algorithmic personalisation systematically reduces viewpoint diversity \parencites[5]{5142478}[29]{5142478}; a professional feed surfacing "thought leadership" through comparable logic would plausibly import bias into hiring and reputational signals (an illustrative inference, not a documented platform finding). Profiling-based advertising documented for Instagram and TikTok \parencites[4]{6941611} could analogously support job ads targeting inferred vulnerability over qualification, compounded by literacy-survey evidence that active users overestimate their manipulation-recognition ability (mean $2.31$/$5$) \parencites[10]{9360054}. Gerr's finding that DSA/AI Act transparency obligations have not removed underlying engagement incentives \parencites[7]{5142478}[28]{5142478} and Chitra et al.'s call for mandatory auditing with conditional liability \parencites[17]{6502223} would apply equally to any professional feed influencing employability, were platform-specific data available. Professional networking is flagged as a future-research priority because such evidence does not yet exist \parencites[5]{5142478}[2244]{6502223}[25]{2977039}. Human-centred fairness metrics \parencites[53]{4905459} are needed because offline accuracy measures miss the systemic harms that professional-opportunity mediation could create.

\subsubsection{Original Heuristic Evaluation: A Single-Evaluator LinkedIn Walkthrough}
The extrapolation above can be partially tested. This report's author walked through LinkedIn's publicly visible desktop and mobile interface and scored observations against the report's own Code A/B/C matrix (Code A interface mechanics, Code B psychological triggers, Code C legal risk tiers). The walkthrough covers only features visible to any logged-in user; it has no access to backend ranking logic, A/B-test variants or internal telemetry, and reflects one evaluator's observations and recollections across sessions. Confidence is lower throughout than for the TikTok and Instagram findings, which rest on a much larger published literature base. The exercise demonstrates the toolkit on a platform with no dedicated dark-pattern literature; it does not claim to close the LinkedIn evidence gap. The result is mixed: some findings confirm the concerning pattern documented elsewhere, others do not.

\begin{table}[H]
\centering
\caption{Heuristic Walkthrough of LinkedIn's Publicly Visible Interface}
\label{tab:linkedin_heuristic_walkthrough}
\begin{tabular}{p{4.6cm}p{4.3cm}p{4.6cm}}
\hline
Feature Observed & Code Classification & Confidence \\
\hline
Continuous feed with no natural end point or session boundary & Code A (interface mechanics) & High (directly observed) \\
Notification and connection-request badges (red indicator) appear on unread items but clear after a single view, with no observed re-escalation or repeat nagging & Code A; low observed Code B salience relative to the variable-reward pattern documented for other platforms & High (directly observed) \\
``Who's viewed your profile'' analytics: some detail available free, fuller detail gated behind a Premium subscription upsell & Code A (economic gatekeeping) / Code C (monetised curiosity, DFA-relevant if repeated at scale) & High (directly observed) \\
Profile-completeness section present; explicit completion nudges were not observed on every session, but 2-3 distinct touchpoints recommended a Premium upgrade to unlock profile-enhancement features & Code A (monetisation-linked) more than a pure Code B goal-gradient effect & Medium (based on this evaluator's recollection across recent sessions, not a single controlled walkthrough) \\
``People you may know'' and connection-request prompts recurring across three distinct surfaces in a single session (feed, dedicated area, notifications); no single-step opt-out located & Code A / Code B & Medium \\
Feed interleaves algorithmically ranked ``suggested'' content among connection updates, not strictly chronological, with no toggle located to switch to a chronological-only view & Code A & High (directly observed, consistent with the ranking-opacity pattern documented for other platforms) \\
Autoplay of video content in the feed & Not observed; videos did not play automatically while scrolling & High (directly observed; noted for completeness, since not every mechanism documented elsewhere generalises to this platform) \\
Subscription downgrade/cancellation flow includes a reason-selection step, but permitted the change without the more severe obstruction patterns documented in the subscription-platform audit elsewhere in this report & Code C, but lower-severity than the audited roach-motel pattern & Medium (single most-recent use recalled by this evaluator, not a fresh systematic test) \\
\hline
\end{tabular}
\end{table}

These observations do not establish prevalence, user harm or regulatory exposure the way the cited TikTok and Instagram literature does. No autoplay was observed; notification badges cleared after a single view without re-escalating; and the cancellation flow examined did not show the severe obstruction documented in the subscription-platform literature. On the other side, the non-chronological feed with no opt-out toggle, connection-nudging across three separate surfaces, and the free/Premium split on profile-view analytics all match the concerning pattern found elsewhere. The Integrity Rule stated at the start of this subsubsection still applies: this walkthrough sits between pure extrapolation and verified platform-specific research. Its mixed result is itself the argument for systematic checking rather than assuming the worst-documented pattern generalises by default, which is precisely what a researcher-administered evaluation of the kind described in Section 6.2 would do properly.

\subsection{Interface Mechanics across Platforms}

\subsubsection{Feed Structures and Scrolling Patterns}
\label{sec:feed_structures_scrolling_patterns} TikTok's default landing page is the personalised "For You" feed \parencites[42]{3195574}, meaning every session begins with algorithm-selected content before any deliberate user choice; micro-duration clips with infinite swipe then eliminate the need to decide again. Li links this architecture to an "auto-pilot mode" associated with weakened social skills and behavioural addiction among college students \parencites[2]{8958037}[6]{8958037}. YouTube's default autoplay follows the same logic, privileging recommended content over manually selected items \parencites[42]{3195574}. Hu connects the resulting feed structure to $22\%$ higher compulsive checking and correlated neurocognitive change \parencites[2]{3303310}[20]{3303310}. The regulatory response follows: the DSA's 2026 TikTok investigation names these features as potential "systematic manipulative design," and Me{\ss}mer and Degeling link default autoplay and personalised-feed landing pages directly to DSA Article 25's dark-pattern prohibition \parencites[6]{8958037}[42]{3195574}. Feed structure also shapes exposure composition beyond session length. Gerr shows engagement-optimised trending reduces viewpoint diversity \parencites[3]{5142478}[29]{5142478}; Instagram's "Explore" ranking increased body-negative imagery exposure by $23\%$ versus chronological feeds \parencites[3]{3303310}. A key audit variable is click-depth to a non-personalised alternative \parencites[42]{3195574}; HITL grading confirms the alternative architecture is viable, embedding oversight without removing the content stream \parencites[17]{6805789}. Infinite scroll and default autoplay are the primary cross-platform inspection targets for the Ethical AI Audit Toolkit \parencites[6]{6295677}[2]{8958037}[2]{3303310}[42]{3195574}.

\subsubsection{Use of Autoplay and Recommendations}
\label{sec:use_autoplay_recommendations} On short-video feeds, the next clip begins before the user acts; stopping requires deliberate interruption, and opting out is the non-default \parencites[2]{8958037}. YouTube's default autoplay receives the same regulatory designation, privileging recommended items over manually selected ones \parencites[42]{3195574}. TikTok compounds the effect by opening on "For You," making both the entry point and every subsequent item algorithm-determined \parencites[42]{3195574}. Algorithmically curated feeds produce $22\%$ higher compulsive checking and $23\%$ higher depressive-symptom incidence alongside reduced prefrontal and grey-matter activity \parencites[2]{3303310}, the pattern Li calls an "algorithmic addiction loop." The DSA's TikTok proceeding treats this combination as actionable \parencites[2]{8958037}[42]{3195574}, and Kostenko's AI Law Model bans autoplay for minors while mandating "detox mode" defaults and visible "no personalisation" switches for teen social networks \parencites[221]{2977039}[223]{2977039}. Indian platform audits document pre-selected defaults operating by the same logic \parencites[4]{6295677}[42]{3195574}. Autoplay and infinite scroll are Code A features; variable rewards and temporal discounting are Code B triggers; age-sensitive design violations are Code C \parencites[2]{8958037}[2]{3303310}[221]{2977039}[6]{6295677}[2]{3303310}[12]{2977039}. The architecture operates identically for adults, who have no comparable regulatory protection.

\subsubsection{Notification Design and Return Incentives}
\label{sec:notification_return_incentives} Push notifications occupy a structurally distinct position from autoplay or infinite scroll: they reach users after they have already left. $47\%$ of users surveyed by Khare et al. reported difficulty cancelling services, and "nagging" notifications with repeated post-refusal prompts generated a documented "wear down effect" steering users toward platform-preferred choices \parencites[8]{6295677}. Urgency-marketing executives describe late-night and high-frequency push campaigns as commercially advantageous, explicitly premised on temporal vulnerability and low self-control windows \parencites[138]{8745065}; combined with stored-payment one-click checkout, these prompts compress deliberation almost to zero. India's CCPA advisory frames this as platform accountability, mandating internal dark-pattern audits and explicit consent over pre-checked defaults, and shifts from complaint-driven enforcement to preventive regulation targeting sludge architecture before it causes harm \parencites[6]{6295677}[12]{6295677}. High-frequency, behaviourally-targeted notifications embedded in obstructive flows are classified here as high-risk manipulative return incentives. Table~\ref{tab:notif_patterns} synthesises these patterns. \begin{table}[H] \centering \caption{Notification related patterns functioning as return incentives} \label{tab:notif_patterns} \begin{tabular}{p{3.2cm}p{5.0cm}p{4.5cm}} \hline Pattern & Mechanism & Documented effects \\ \hline Nagging notifications & Repeated prompts after refusal, often via app alerts & Wear down effect, decision fatigue, coerced consent \parencites[8]{6295677} \\ Late night urgency messages & Time based push for expiring discounts or offers & Executives perceive higher responsiveness in "low control" windows \parencites[138]{8745065} \\ Subscription renewal prompts inside obstructive flows & Notifications point to renewal while exit paths remain complex & Forced continuation of paid services, recurring unwanted payments \parencites[8]{6295677} \\ "Dark pattern free" self declarations & Internal audits and voluntary declarations without effective external verification & $81\%$ of platforms self declared dark pattern free despite persistent complaints, 28\% resolution rate \parencites[12]{6295677} \\ \hline \end{tabular} \end{table} Bansal's synthesis of algorithmic advertising architecture shows notifications bind users through three stacked layers: a "cognitive trigger layer" deploying scarcity cues; an "algorithmic personalization layer" tailoring timing to behavioural histories; and a "temporal and friction reduction layer" amplifying urgency while shrinking steps to purchase \parencites[139]{8745065}. Exploitative impulse buying is an architectural outcome of this stack. TikTok's infinite scroll, autoplay and personalised feeds are already under DSA investigation as "systematic manipulative design" \parencites[2]{8958037}; behaviourally tuned notifications re-routing users back into those feeds replicate the same pattern in a second channel. Two configurations reach Code C: notifications embedded in obstructive exit flows or subscription traps, and those combining temporal-vulnerability assumptions with high-urgency messaging documented to compress deliberation \parencites[8]{6295677}[138]{8745065}. The India CCPA's mandate for preventive auditing reflects the shared conclusion that current self-declaration is insufficient.

\subsection{Psychological Trigger Mapping by Platform}

\subsubsection{Variable Rewards and Intermittent Reinforcement}
\label{sec:variable_rewards_intermittent} Short-video platforms engineer intermittent reinforcement into default feed architecture. Li describes neural-network recommenders paired with swipe-based feeds that trigger dopamine release during the expectancy moment before the next clip, producing "auto-pilot" scrolling associated with diminished academic concentration, heightened anxiety and depression among college students \parencites[2]{8958037}[3]{8958037}[6]{8958037}. Hu extends the evidence to neurocognitive change: $22\%$ higher compulsive-checking behaviour, $23\%$ higher incidence of depressive symptoms versus chronological content, and reduced prefrontal cortex activity and grey-matter density in impulse-control regions \parencites[2]{3303310}[20]{3303310}. Table~\ref{tab:vr_patterns} maps these triggers to working Code B categories. \begin{table}[H] \centering \caption{Variable reward patterns and intermittent reinforcement across platforms (working Code B mapping)} \label{tab:vr_patterns} \begin{tabular}{p{3.2cm}p{4.8cm}p{4.8cm}} \hline Platform type & Variable reward mechanic & Documented effects \\ \hline Short video feeds & Swipe to next clip with unpredictable quality; infinite scroll; default auto play & "Auto pilot" scrolling, behavioural addiction, impaired study performance in students \parencites[2]{8958037}[6]{8958037} \\ Algorithmic social feeds & Engagement ranking that privileges emotionally charged content; irregular delivery of highly salient posts & $22\%$ higher compulsive checking; 23\% higher depressive symptoms; reduced prefrontal activity \parencites[2]{3303310}[20]{3303310} \\ E commerce flows & Countdown timers and "only two left" alerts that sometimes expire and sometimes do not & Conversion uplift reported by executives; decision fatigue and regret in consumers \parencites[136]{8745065}[8]{6295677} \\ \hline \end{tabular} \end{table} Countdown timers and "only two left" flags work precisely because their inconsistency is itself a variable reward \parencites[136]{8745065}[7]{6295677}: users learn that delay sometimes costs them, reinforcing impulsive action without stable informational value. The EU DSA proceeding against TikTok cites infinite scrolling, autoplay and push notifications as possible "systematic manipulative design" \parencites[2]{8958037}. Kostenko's AI Law Model goes further, prohibiting systems that cause digital addiction and establishing an inalienable right to information environments free from chronic anxiety-inducing flows \parencites[61--62]{2977039}[119]{2977039}. Hu's interdisciplinary matrix demonstrates audit measurability: neurocognitive metrics map to algorithmic feature vectors via the Algorithmic Impact Assessment Scale (AIAS-5), quantifying emotional volatility, sleep disruption and compulsive usage \parencites[19]{3303310}[20]{3303310}. Feeds combining infinite scroll, default autoplay and irregular high-arousal content are classified Code C in this report's working taxonomy, particularly for users who are confident in digital skills yet score poorly at recognising manipulation \parencites[10]{9360054}[2]{3303310}[2]{8958037}. Recommender systems sustain this engagement via collaborative filtering and network-based models \parencites[3]{7422270} that scale efficiently on large sparse interaction datasets, making the reinforcement loop technically resilient where regulation has not yet reached.

\subsubsection{Social Validation, Status Cues and Fear of Missing Out}
\label{sec:social_validation_status_fomo} On Douyin and Xiaohongshu, visible engagement metrics are unreliable signals. Wang reports that likes, follower counts and reviews are routinely distorted by fake accounts and paid engagement \parencites[13]{7160486}; a 2023 China Consumers Association report found nearly $30\%$ of online shoppers had encountered falsified reviews, converting aggregated trust signals into manipulation instruments. Influencer and KOL endorsements compound the distortion because personal credibility and paid promotion are deliberately fused, even where SAMR sponsorship-disclosure rules apply \parencites[13]{7160486}. Algorithmic curation then amplifies whatever is already popular. Douyin's viral challenges and livestreamed sales create collective-buying moments where social pressure drives impulsive spending and emotional fatigue \parencites[13]{7160486}, and Instagram's "Explore" page amplified body-negative imagery by $23\%$ versus chronological feeds, correlating with statistically significant self-esteem declines among adolescent women of color \parencites[2]{3303310}[3]{3303310}. Table~\ref{tab:status_fomo} summarises these triggers in the working coding matrix. \begin{table}[H] \centering \caption{Working Code B mapping: social validation, status cues and FOMO} \label{tab:status_fomo} \begin{tabular}{p{3cm}p{4.5cm}p{4.5cm}} \hline Pattern & Mechanism & Documented effects \\ \hline Inflated metrics & Fake likes, followers, paid engagement & Illusion of popularity, weakened trust in reviews and ratings \\ Undisclosed sponsorship & Influencer posts that hide commercial nature & Exploitation of emotional credibility, confusion between advertising and opinion \\ Algorithmic herd reinforcement & Preferential promotion of already popular content & Pressure to conform to trends, FOMO driven participation in viral challenges and live sales \parencites[13]{7160486} \\ Body and identity sensitive feeds & Prioritisation of high arousal or body negative imagery & Increased body dissatisfaction and lower self esteem among specific groups \parencites[3]{3303310} \\ \hline \end{tabular} \end{table} Status-seeking shapes purchasing behaviour across these same platforms. Kishore et al. find university students' marketplace choices dominated by price, urgency and reviews; ethical considerations remain secondary, and unverifiable social proof triggers scepticism that widens the intention-behaviour gap \parencites[2]{4245543}[4]{4245543}. Yin reports TikTok hashtag challenges make participants $62\%$ more likely to log in daily \parencites[3]{7154741}, linking routine engagement to status contests where missing a challenge risks social invisibility. Distorted social proof, herd-reinforcing algorithms and FOMO mechanics are treated as high-priority Ethical AI Audit Toolkit items on the strength of overlapping evidence from Chinese e-commerce, social-media research and student-behaviour studies \parencites[13]{7160486}[2]{3303310}[2]{4245543}[3]{7154741}.

\subsubsection{Emotional High-Arousal and Outrage Content}
\label{sec:emotional_high_arousal_outrage} Chitra et al. show moral-emotional language in tweets raises retweet probability by roughly $20\%$ per additional moral-emotional word \parencites[7]{6502223}, and content evoking awe, anger or anxiety is significantly more shareable than sad content. A ranking stage trained on click-through, watch time, completion and comment signals therefore selects for outrage-maximising content by construction, with no term for epistemic quality in the loss function. Hu documents the measurable result: Instagram's "Explore" page amplified body-negative imagery by $23\%$ relative to chronological feeds, correlating with a statistically significant self-esteem decline among adolescent women of color ($\beta = -0.41$, $p < 0.05$), and algorithmically prioritised content exhibits $31\%$ higher negative affect than baseline \parencites[3]{3303310}[2]{3303310}. Yin documents both directions of viral mobilisation: Black Lives Matter catalysed global protest, while COVID-19 misinformation spread faster than official retraction as algorithms promoted "racy" high-interaction posts, with creators reporting harassment and burnout and viewers reporting distress from repeated exposure \parencites[2]{7154741}. Table~\ref{tab:high_arousal_mapping} aligns these with the working coding matrix. \begin{table}[H] \centering \caption{High arousal and outrage content: working Code B mapping} \label{tab:high_arousal_mapping} \begin{tabular}{p{3.2cm}p{4.6cm}p{4.6cm}} \hline Coding slot & Mechanism & Empirical or conceptual anchor \\ \hline B1: Moral emotional amplification & Moral emotional wording boosts sharing, driving outrage selection & Moral emotional words increase retweets by $\approx 20\%$ per word \parencites[7]{6502223} \\ B2: Negative affect bias & Ranking favours high negative affect content & Algorithmic feeds show $31\%$ higher negative affect, $23\%$ more body negative imagery \parencites[3]{3303310} \\ B3: Viral echo chamber reinforcement & Viral, racy posts amplified into cohesive cliques & Viral misinformation and doctored videos deepen echo chambers and division \\ B4: Emotional fatigue and distress & Repeated exposure to polarising material harms creators and viewers & Harassment, burnout, information fatigue and emotional distress reported \parencites[2]{7154741} \\ \hline \end{tabular} \end{table} Candidate-generation and ranking stages optimise multi-task objectives weighting click-through, completion, watch time, share and comment rate by commercial value, with no wellbeing term in the loss function. YouTube's "radicalization pipeline" follows from this: a reinforcement-learning agent that discovers escalating content intensity sustains attention will gradually steer users toward more extreme material \parencites[7]{6502223}. Chitra et al. frame algorithmic manipulation of political information as threatening freedom of expression and electoral rights \parencites[11]{6502223}; digital-constitutionalist proposals call for ranking transparency, audit trails, independent oversight and non-discrimination in political-information distribution. Emotional high-arousal and outrage content is a predictable outcome of engagement-oriented optimisation \parencites[7]{6502223}[3]{3303310}[2]{7154741}, making these patterns central diagnostic signals for any audit toolkit aligned with EU digital constitutionalism \parencites[7]{6502223}[11]{6502223}[3]{3303310}.

\subsection{Audit Table and Findings}

\subsubsection{Consolidated Cross-Platform Risk Matrix}
\label{sec:audit_table_findings} Sections 4.1 through 4.3 produce four platform-specific tables (Tables~\ref{tab:notif_patterns}, \ref{tab:vr_patterns}, \ref{tab:status_fomo} and \ref{tab:high_arousal_mapping}). Table~\ref{tab:cross_platform_matrix} consolidates these into a single working matrix, mapping each platform type against its dominant Code A/B feature, Code C risk tier and the SDT indicator most directly degraded. LinkedIn is absent: per the Integrity Rule applied in Section 4.1.3, no source names the platform specifically, so it is scoped out rather than filled with extrapolated figures. \begin{table}[H] \centering \caption{Cross-platform audit matrix: features, DFA-adjacent risk tier, and SDT degradation indicator} \label{tab:cross_platform_matrix} \begin{tabular}{p{2.6cm}p{3.4cm}p{2.1cm}p{4.4cm}} \hline Platform type & Dominant Code A/B feature & Code C tier & SDT indicator degraded \\ \hline Short-video (TikTok-type) & Infinite scroll, autoplay, variable reward \parencites[2]{8958037}[2]{3303310} & High/critical & Autonomy (auto-pilot scrolling); competence (impaired concentration) \\ Image/story (Instagram-type) & Explore-style ranking, social proof, body-negative amplification \parencites[3]{3303310} & High & Relatedness (distorted social comparison); competence (self-esteem decline) \\ Professional networking (illustrative only) & Engagement ranking, profiling-based ads (extrapolated, not platform-verified) \parencites[2]{6502223}[4]{6941611} & Not yet assessable & Not yet assessable, flagged for future primary research \\ Cross-platform notification layer & Nagging, obstructive cancellation, temporal-vulnerability targeting \parencites[8]{6295677}[138]{8745065} & High & Autonomy (compressed deliberation) \\ Cross-platform outrage/virality layer & Moral-emotional amplification, negative-affect bias \parencites[7]{6502223}[3]{3303310} & High & Relatedness (echo-chamber reinforcement); autonomy (opinion manipulation) \\ \hline \end{tabular} \end{table} Two findings emerge from reading the matrix as a whole. Every high-tier row combines a Code A interface mechanic with a Code B psychological trigger; none of the documented harms traces to either category in isolation. The "not yet assessable" professional-networking row is an indicator that audit coverage tracks the literature's own gaps, not the audit team's priorities.

\subsection{Synthesis of Core Vulnerabilities}

\subsubsection{Cross-Cutting Patterns Behind the Platform Findings}
\label{sec:synthesis_core_vulnerabilities} Across every platform examined in Section 4, the same structural weaknesses reappear despite changes in surface architecture.

Compressed deliberation runs through every product category. Countdown-style urgency, nagging notifications and obstructive cancellation flows all shorten the window for System 2 reflection \parencites[8]{6295677}[138]{8745065}[7]{6295677}, regardless of whether the surface is a short-video feed, a subscription paywall or a notification tray.

Reward unpredictability is equally portable. Infinite scroll, variable-reward swipe mechanics and algorithmically amplified high-arousal content share identical intermittent-reinforcement logic \parencites[2]{8958037}[2]{3303310}: swap short clips for outrage-bait posts and the addictive mechanism transfers intact.

What ties these together is an asymmetry of awareness. Users consistently rate themselves digitally competent while scoring poorly on recognising manipulation (mean $2.31/5$) \parencites[31]{9360054}[7]{1456534}. Disclosure-only remedies such as transparency reports systematically under-protect the population the DSA's non-profiled feed requirement is meant to reach.

These three vulnerabilities co-occur and reinforce one another. A platform combining compressed deliberation, unpredictable rewards and low user awareness is not accumulating three separate risks; confirming any one should prompt immediate checking of the other two. That co-occurrence is the core finding this report's Ethical AI Audit Toolkit (Section 5) is built to act on.

\section{Practical Recommendations and Pro-Social Pivot}

\subsection{Conceptual Design of the Audit Toolkit}

\subsubsection{Objectives, Users and Governance of the Toolkit}
\label{sec:toolkit_objectives_users_governance} Three linked objectives organise audit design here: legal and ethical duties must become operational control points rather than aspirational statements; artefacts generated must be inspectable by external actors beyond internal teams alone; and controls must embed into standing governance as permanent checkpoints.

Catapang's triple-gate architecture demonstrates what practical, enforced control looks like \parencites[24]{8367105}. Governance, Metric and Eco gates block or throttle releases at each pipeline stage, mapped to AI Act obligations and NIST AI RMF functions, and generate logs as audit artefacts. The toolkit adapts this model to prevent feed-ranking releases that worsen addiction or fairness metrics beyond predefined thresholds, thereby carrying out Kostenko's mandatory ethical-admissibility procedures \parencites[9]{2977039}: notification, explainability, and human control over critical decisions. Gate logic is the mechanism; logged artefacts are the proof.

Envisioned user groups differ in what they require. Product and engineering teams need CI/CD-integrated gate checks aligned to DSA non-profiled-option obligations \parencites[12]{6502223}. Risk and compliance functions need documentation and scoring outputs comparable to proposed algorithmic impact assessments for credit AI \parencites[15]{5442231} and governance-matrix scoring sheets \parencites[59]{2702055}. External regulators face documented systemic gaps in auditing capacity \parencites[9]{2977039}; Bahangulu and Owusu Berko argue that businesses must adopt external oversight committees precisely because regulators struggle with black-box models \parencites[8]{7461574}[9]{7461574}. The toolkit bridges both groups through observable outputs and logged gate decisions.

Combining internal ethics boards, external oversight and statutory frameworks yields measurably better bias reduction and appeal resolution than self-regulation alone \parencites[2]{0375818}. Pastoral Data Trust governance models prevent "digital colonialism" in fragile mental-health deployments by barring any single party from unilaterally controlling data access \parencites[11]{5419804}[13]{5419804}. This thesis situates toolkit governance co-regulatorily: internal ethics committees manage gate configurations while independent auditors and national registers access gate artefacts and risk assessments. HITL grading's academic oversight committees provide a working precedent for rubric standardisation and recalibration \parencites[21]{6805789}. Most ethical failures trace back to the control architecture rather than the stated principle \parencites[24]{8367105}[9]{7461574}[12]{6502223}[26]{2977039}; an enforced gate that blocks the release matters more than a policy document that prohibits the feature.

\subsubsection{Audit Typologies Relevant for Recommender Systems}
\label{sec:audit_typologies_recommenders} Four complementary typologies from adjacent domains inform audit design here.

The first is the compliance-versus-scenario distinction Me{\ss}mer and Degeling draw for DSA Article 37 work. Compliance audits verify documented risk assessments. Scenario-based risk audits apply the ISO five-step process to socio-technical testing of ranking behaviour under defined conditions such as self-harm content exposure \parencites[21]{3195574}. Empirical archetype-account studies show that such scenario tests can reveal algorithmic saturation toward extreme material within approximately five days of passive viewing \parencites[4]{2580223}. Short behavioural probes surface harmful dynamics that documentation-only reviews miss.

The second typology is the algorithmic impact assessment proposed by Sebastian for high-risk credit-AI systems: an environmental-impact-review equivalent testing fairness, accuracy and systemic effects both pre- and post-deployment \parencites[15]{5442231}. Engagement-optimising feeds meet this threshold. Periodic assessment is warranted well beyond content legality alone.

Third, sector-specific governance matrices provide a scoring template. Mutawa et al.'s Algorithmic Governance Assessment Matrix for zakat institutions shows that transparency scores correlate with documentation clarity, and AI readiness scores with eligibility accuracy \parencites[4]{2702055}[8]{2702055}. The same matrix logic transfers to recommender platforms.

Fourth, Kostenko's AI Law Model separates design-time ethical admissibility tests from ongoing risk-tiered supervisory audits by accredited auditors \parencites[8]{2977039}[9]{2977039}. Admissibility tests cover notification, explainability and human control over critical decisions. Auditors may reclassify risk based on changing operational metrics; reclassification is not a sanction but the intended supervisory loop functioning correctly.

Financial-AI lifecycle validation pipelines add a further process pattern: adversarial pre-deployment bias testing, a self-audit loop logging SHAP attributions to a write-once ledger, bi-annual third-party review, and automatic circuit-breaker fallback when drift thresholds are breached \parencites[7]{8907572}. HITL grading demonstrates the same principle in educational AI. There, operational audits are inseparable from daily workflow, and oversight committees analyse override logs and recalibrate models continuously \parencites[21]{6805789}. This report's toolkit aligns its modules to these four typologies, subject throughout to the black-box constraint introduced in Section 5.1.3.

\subsubsection{Alignment with Regulatory and Ethical Requirements}
\label{sec:alignment_regulatory_ethical} Toolkit choices must trace to hard-law and model-law obligations. Any audit item that cannot be grounded in a binding duty or authoritative model standard risks becoming a cosmetic layer dismissed during enforcement review.

The AI Act's unacceptable-risk tier and Kostenko's ban on behaviour-manipulating systems jointly anchor Code Category C flags for addictive feeds and minor profiling \parencites[2248]{6502223}[10]{2977039}[61--62]{2977039}. Platform-facing DSA duties form a second anchor: the non-profiled-feed requirement, annual independent audits, and Article 25's ban on interfaces that deceive or "substantially impair" free decisions translate directly into observable audit checks on default personalisation and exit-flow friction \parencites[2248]{6502223}[20]{6941611}[6]{6295677}, requiring no internal ranking-code access.

Kostenko's model law extends these obligations into finer-grained checklist items: mandatory notification, explainability and human-control procedures are each individually required \parencites[8]{2977039}, with technical passports and event logging additionally mandated for high-risk systems \parencites[9]{2977039}[19.5]{2977039}. Tamper-evident logs and human-override evidence are compliance deliverables.

ISO/IEC risk-management, data-quality and bias-assessment standards \parencites[4]{8589961} support AI Act fairness requirements even where technically voluntary. Adopted as internal metrics crosswalked to AI Act criteria, they mirror Catapang's triple-gate regulatory mapping \parencites[23]{8367105}.

HITL grading and AI-proctoring research supply concrete communication design patterns (explicit AI/human tagging, dual logging, privacy-by-design consent \parencites{6805789}[20]{6805789}[9]{2055956}) informing audit items on user-facing disclosure and appeal accessibility. Within the evidentiary limits of secondary qualitative research, the toolkit remains congruent with AI Act, DSA and AI Law Model expectations \parencites[2248]{6502223}[20]{6941611}[8]{2977039}[9]{2055956}.

\subsection{Audit Preparation and Scoping Modules}

\subsubsection{Defining System Boundaries and Use Cases}
\label{sec:defining_system_boundaries_use_cases} Scoping an audit too narrowly misses the manipulation vectors that matter most. This toolkit treats recommenders as black-box services whose internal weights cannot be inspected; assessment relies on observable interface behaviour, regulatory records and secondary studies \parencites[7]{6295677}[2]{8958037}[2]{3303310}.

Boundaries are articulated along three axes. Functional scope draws on Mutawa et al.'s governance matrix to treat the recommender as comprising ranking logic, data capture, user-facing interfaces and surrounding governance mechanisms \parencites{2702055}; for SOCIO SPHERE AI, this must include sentiment analysis and AI-assisted replies alongside ranking proper \parencites{3413731}. Institutional scope follows Kostenko's distinction between AI providers and operators \parencites[20]{2977039}, and extends to include internal ethics boards and external auditors \parencites[2]{0375818}{2702055}. Jurisdictional scope follows Kostenko's technological-sovereignty provisions \parencites[108]{2977039}, alongside AI Act and DSA transparency duties and Chitra et al.'s call for mandatory auditing of recommenders influencing political discourse \parencites[16]{6502223}. Each audit must specify which stack components fall under EU-style duties and which under model-law expectations.

Table~\ref{tab:use_cases_scope} illustrates three archetypal use cases recurring in later chapters. \begin{table}[H] \centering \caption{Illustrative use case archetypes and boundary focus (working taxonomy)} \label{tab:use_cases_scope} \begin{tabular}{p{3.2cm}p{4.4cm}p{4.4cm}} \hline Use case archetype & Boundary focus & Primary sources informing scope \\ \hline Short form video feed for students & Infinite scroll, auto play, push notifications, personalisation pipeline, child data processing, operator governance & Algorithmic addiction loop, student harms, DSA proceeding against TikTok, age appropriate design and child data processing duties \parencites[2]{8958037}[6]{8958037}[2]{3303310}[20]{2977039} \\ Image centric social discovery & Explore style ranking, social proof widgets, behavioural advertising modules, bias and body image impacts & Amplification of body negative imagery, negative affect bias, social proof manipulation and its impact on trust and identity \parencites[3]{3303310}[13]{7160486} \\ AI augmented social platform (SOCIO SPHERE AI) & AI recommendation layer, sentiment analysis, AI assisted replies, moderation AI, admin analytics dashboards & Three tier plus AI/ML architecture, AI generated captions and hashtag suggestions, AI based moderation, engagement analytics \parencites{3413731} \\ \hline \end{tabular} \end{table} Each archetype defines the "system under audit" differently. The student short-video case centres child profiling and predictive "life path" scoring under enhanced-protection provisions \parencites[20]{2977039}[62]{2977039}; SOCIO SPHERE AI must include sentiment analysis and AI-assisted replies even where models are shared across surfaces, meaning a shared recommendation model between main feed and notifications layer requires auditing across both \parencites{3413731}.

Use cases also fix expected artefacts. Following Mutawa et al.'s governance-matrix logic \parencites{2702055}, each use case specifies expected passports, audit logs and impact assessments per Kostenko's transparency chapter \parencites[68]{2977039}[69]{2977039}, while explicitly excluding inaccessible training corpora and live A/B results. The toolkit evaluates declared objectives, interface mechanics and published outputs \parencites[68]{2977039}[21]{3195574}[2]{0375818}.

\subsubsection{Stakeholder Mapping and Risk Prioritization}
\label{sec:stakeholder_mapping_risk_prioritization} Knowing who carries legal duties is the foundation of risk prioritisation. Kostenko distinguishes developers, operators, regulators and affected individuals, all bound by ethical-responsibility duties \parencites[61--62]{2977039}. The Algorithmic Impact Navigator similarly positions deploying organisations as "the first line of accountability" \parencites[2]{2742041}[3]{2742041}.

Four stakeholder clusters are central for audit design. Platform organisations include leadership commissioning impact evaluations, product and engineering teams building the ranked feed, and ethics committees responsible for gate configurations \parencites[2]{2742041}[61--62]{2977039}. External oversight actors span EU regulators, national authorities, and the accredited auditors central to co-regulatory models \parencites[10]{2977039}[2]{0375818}. Users and affected communities include heavy social-media users with low manipulation awareness \parencites[10]{9360054}, students experiencing short-video addiction \parencites[2]{8958037}, and communities subject to AI-mediated mental-health interventions \parencites[13]{5419804}. Each cluster carries different legal standing and different information needs. Epistemic stakeholders (researchers and civil society) rely on data access \parencites[21]{3195574} to study systemic risk and may trigger regulatory enforcement. Where their access rights under DSA Article 40 analogues go unused or unenforced, that absence is itself a risk signal worth tracking.

Risk prioritisation follows the AIN's screening logic around sociopolitical context, opacity and critical-infrastructure integration \parencites[3]{2742041}. Engagement-optimised feeds heavily used by students rank high: vulnerability is elevated, ranking is opaque, and integration into daily routine is pervasive. Chitra et al.'s argument for mandatory researcher data access to very-large-platform recommenders affecting civic discourse reinforces this prioritisation under the DSA \parencites[2248]{6502223}. Table~\ref{tab:stakeholder_risk_priorities} outlines a working mapping. \begin{table}[H] \centering \caption{Stakeholder centric risk priorities for social recommender audits (working taxonomy)} \label{tab:stakeholder_risk_priorities} \begin{tabular}{p{3cm}p{4.2cm}p{4.5cm}} \hline Stakeholder cluster & Primary concern & High priority risk cues \\ \hline Platform leadership and boards & Legal compliance, reputation, business continuity & Presence of addictive design patterns under DSA scrutiny; absence of internal ethical review and AI Ethics Impact Reports \parencites[2]{8958037}[61--62]{2977039} \\ Engineering and product teams & Technical performance, shipping velocity & Lack of gating controls in CI/CD for fairness, transparency and sustainability; missing logs needed for conformity assessment \parencites[24]{8367105}[19.5]{2977039} \\ Regulators and auditors & Systemic risk, fundamental rights & Inability to obtain technical dossiers and logs; absence of algorithmic impact assessments and standardised fairness metrics \parencites[19]{3195574}[15]{5442231} \\ Users and communities & Autonomy, psycho emotional safety, fairness & Infinite scroll and auto play combined with personalised feeds; documented mental health or discrimination harms \parencites[2]{8958037}[2]{3303310}{4245543} \\ Researchers and civil society & Transparency, contestation & Restricted data access under DSA Article 40 analogues; opaque terms and conditions on content and risk management \parencites[19]{3195574} \\ \hline \end{tabular} \end{table} Governance model also modulates risk level. Hybrid governance yields greater bias reduction and audit compliance than self-regulation alone \parencites[2]{0375818}; zakat institutions with higher AI readiness and transparency scores achieve better eligibility accuracy \parencites{2702055}. Platforms under co-regulatory or statutory regimes with documented ethics councils therefore rank lower in structural risk than those relying on voluntary self-assessment.

One methodological limit applies throughout. Dark-pattern prevalence is assessed via secondary institutional reports \parencites[7]{6295677}; engagement-related harms are inferred from controlled experiments, not internal telemetry \parencites[2]{3303310}. Stakeholder mapping and risk prioritisation derive from documentary and empirical sources describing behaviours external to proprietary systems \parencites[10]{2977039}[21]{3195574}[9]{7461574}. The toolkit is a qualitative documentation-centric instrument, a deliberate scope boundary.

\subsubsection{Documentation and Data Access Pre-Checks}
\label{sec:doc_data_prechecks} Before any scenario-based or fairness assessment begins, auditors must establish a simpler baseline: does the required documentation exist? Regulators expect specific artefacts, and their absence signals problems beyond any particular audit finding.

Kostenko's TSAI regime provides a detailed checklist. High-risk systems require a unique identifier and an algorithmic passport specifying owner, risk class, model architecture, data sources, oversight channels and event-logging retention; absence of an equivalent dossier signals serious traceability gaps. Data-access duties require owners to preserve and provide event logs, with failure triggering fines or access suspension \parencites[244]{2977039}. DSA Article 40 positions vetted researcher data access as central to systemic risk audits \parencites[32]{3195574}. Pre-checks must verify that log retention meets minimum periods (at least $12$ months under the AI Law Model) and that internal procedures exist for supplying materials on request.

Logging quality matters as much as logging existence. Kostenko requires logs coupled to appeal mechanisms suspending disputed decisions pending review \parencites[191]{2977039}. HITL grading provides a practical compliance model: every decision is time-stamped and linked to reviewer identifiers, forming an inspectable override-pattern audit trail \parencites[16]{6805789}[17]{6805789}. Pre-checks confirm whether platforms maintain model cards and change logs per Kostenko's model risk-management provisions \parencites[191]{2977039}; a gap here indicates governance intention without the evidentiary infrastructure compliance requires.

A structured pre-check grid recording the existence and quality of passports, event logs, appeal channels and audit histories \parencites[244]{2977039}[15]{5442231}[59]{2702055} flags missing or inadequate artefacts as Code Category C: "Critical Traceability Failure." Under AI Law Model expectations, qualification for high-risk deployment is conditional on evidence \parencites[244]{2977039}.

\subsection{Interface and Choice Architecture Audit Checklists}

\subsubsection{Assessment of Feed Mechanics and Session Design}
\label{sec:assessment_feed_mechanics_session_design} Feed mechanics sit at the heart of algorithmic addiction. Infinite scroll, default autoplay and personalised landing views collectively determine how easily users slip into extended sessions and how difficult those sessions are to interrupt. Short-video platforms exemplify this via a closed loop of "personalized $\rightarrow$ need satisfaction $\rightarrow$ behavioural data accumulation $\rightarrow$ re-push" \parencites[2]{8958037}[3]{8958037}; Hu reports $22\%$ higher compulsive checking and reduced grey-matter density on algorithmically curated feeds \parencites[2]{3303310}[20]{3303310}.

Four explicit audit checks follow: default entry point, presence of infinite scroll, autoplay behaviour, and accessibility of a non-personalised option. TikTok's default "For You" launch and YouTube's default autoplay already face DSA scrutiny as potential "systematic manipulative design" \parencites[42]{3195574}[2]{8958037}. Configurations combining always-personalised launch, infinite scroll, default autoplay, and no accessible non-profiled option are coded Code Category C: "addiction-enabling feed mechanics." Douyin averages 21+ daily sessions, with $56.9\%$ of users exceeding 60 minutes, particularly among adolescents with weaker self-control \parencites[3]{8958037}, and the $23\%$ higher depressive-symptom incidence linked to hyper-personalised feeds reinforces this classification \parencites[20]{3303310}.

Mitigating structures must also be documented. HITL grading shows that continuous streams can accommodate autonomy: AI scores remain subject to educator validation, with $87\%$ accepted with only minor edits \parencites{6805789}[19]{6805789}. Auditors should ask equivalently whether feeds expose "why am I seeing this" overlays and whether humans can intervene in ranking with logged effect, since logged intervention is what distinguishes auditable control from stated intention.

\subsubsection{Evaluation of Friction, Defaults and Opt-Out Channels}
\label{sec:evaluation_friction_defaults_optout} Meaningful control depends not only on whether settings exist, but on whether users can reach and activate them without effort or penalty. Dark-pattern research in India finds interface interference and obstruction on roughly $70\%$/$63\%$ of major platforms; $55\%$ of consumers report unexpected charges and only $39\%$ recognise the term "dark patterns" \parencites[2]{6295677}[10]{9360054}. Friction and defaults erode autonomy precisely where digital literacy is weakest.

DMA choice-screen research shows that interoperability mandates reduce switching friction only with vigilant enforcement against manipulative recreation of lock-in \parencites[10]{1070605}; calls for bright-line prohibitions signal that complaint-driven enforcement is structurally inadequate against sludge architectures \parencites[13]{6295677}[14]{6295677}. Asymmetric defaults and obstructive flows are coded Category C: "high-risk friction and default configurations," checked across three loci.

On consent: India's CCPA advisory requires "explicit user action" over pre-checked defaults. The gap between $81\%$ of platforms self-declaring "dark pattern free" and complaint resolution sitting at $28\%$ warrants Code C flagging regardless of formal notices \parencites[12]{6295677}.

On subscriptions: roach-motel patterns combining easy entry with obstructive exit \parencites[4]{6295677} are assessed against the AI Law Model's ban on interfaces that make refusal difficult \parencites[207]{2977039}; account deletion or feed deactivation must match sign-up in step count.

On notifications: Khare et al.'s "nagging" pattern \parencites[14]{6295677} is assessed against Kostenko's age-appropriate-design limits on push notifications for minors \parencites[223]{2977039}; refusing notifications must not degrade core service access. These asymmetric defaults and obstructed controls are high-salience manipulation indicators warranting prioritisation alongside feed mechanics \parencites[2]{6295677}[31]{9360054}[223]{2977039}.

\subsubsection{Classification of Nudging Practices along an Ethics Spectrum}
\label{sec:classification_nudging_spectrum} Not every nudge is a dark pattern. Interface elements that guide users become problematic only when guidance secures compliance by degrading comprehension, autonomy or reversibility.

Liu and Keshavarz distinguish online choice architecture generally from dark commercial patterns, which "steer, deceive, pressure, or obstruct consumers" \parencites[3]{6087241}. The ethics spectrum has two dimensions: the extent to which a design supports preference-consistent choice, and the presence or absence of manipulation indicators; both dimensions must be scored.

At the acceptable end, "assistive nudges" (neutral rankings, symmetric accept/reject consent dialogs) preserve informed choice while still persuading \parencites[33]{6087241}, coded low risk when paired with transparent CSR governance and regular audits \parencites[16]{8200246}. Mid-spectrum "borderline nudges" raise short-term conversion but risk eroding trust; cookie notices making acceptance easier than refusal, found across large-scale crawls of roughly $11{,}000$ shopping sites, exemplify this band \parencites[33]{6087241} and require remediation and complaint/churn monitoring rather than immediate prohibition.

At the unacceptable end sit "coercive sludge patterns," combining exploitative nudging with artificial friction. Khare et al. document roughly $87\%$ of major Indian platforms deploying at least one such pattern: drip pricing, false urgency, roach motels, forced action. Each exploits System 1 heuristics while blocking System 2 reflection through complexity \parencites[7]{6295677}[4]{6295677}, warranting Code Category C: "Prohibited or Critical Risk Nudges," consistent with EU Digital Fairness discussions and AI law models seeking to eliminate manipulative design \parencites[390]{7896734}[61--62]{2977039}.

Table~\ref{tab:nudging_spectrum} applies this ethics spectrum inside audits: \begin{table}[H] \centering \caption{Working ethics spectrum for nudging practices in interface and choice architecture} \label{tab:nudging_spectrum} \begin{tabular}{p{3cm}p{4.2cm}p{4.2cm}} \hline Category (working label) & Typical design attributes & Indicative sources / risk signals \\ \hline Assistive nudges (low risk) & Symmetric effort for accept / decline, clear total costs, visible lower priced alternatives, easy reversal & Preserves informed, preference consistent choice; aligns with ethical digital marketing and CSR based AI governance \parencites[33]{6087241}[16]{8200246} \\ Borderline nudges (medium risk) & Mild interface asymmetry, moderate highlighting of defaults, incomplete but non deceptive urgency messages & Asymmetric cookie banners where refusal is possible but less salient; potential conversion trust trade off \parencites[33]{6087241} \\ Coercive sludge patterns (Code C) & Hidden fees, false scarcity, confirm shaming, roach motels, obstructive cancellation, forced action & Documented as structural features on $\approx 87\%$ of major platforms; linked to autonomy erosion and regulatory calls for bright line bans and UX audits \parencites[7]{6295677}[14]{6295677}[390]{7896734} \\ \hline \end{tabular} \end{table} The anticipated Digital Fairness Act targets "addictive product designs and unfair personalisation," complementing DSA and DMA bans on interface pressure and the AI Act's prohibition of subliminal manipulation \parencites[390]{7896734}[21]{6941611}; Kostenko's confidentiality principle additionally forbids profiling-driven influence without explicit consent \parencites[44]{2977039}. Taken as a set, these instruments provide legal backing for treating Code C patterns as incompatible with emerging EU digital fairness norms across jurisdictions. Triangulating documentary sources with large-scale user-generated evidence and a transparent scoring rubric \parencites[5]{4499345} helps auditors distinguish intentional manipulation from incidental friction.

\subsection{Pro-Social Choice Architecture Patterns}

\subsubsection{Finite Session Design and Break Prompts}
\label{sec:finite_session_break_prompts} Personalised streams produce \parencites[2]{3303310}[20]{3303310} $22\%$ higher compulsive-checking rates and $23\%$ higher depressive-symptom incidence than chronological alternatives, with neuroimaging showing reduced prefrontal activity in heavy users. The DSA proceeding against TikTok names \parencites[2]{8958037} infinite scrolling, autoplay and push notifications as potential ``systematic manipulative design'', and Kostenko's AI Law Model bans \parencites[223]{2977039}[221]{2977039} these mechanics for minors while mandating a ``detox mode'' default with visible ``no personalisation'' switches.

A compliant finite-session design rests on three elements. Soft time-based prompts triggered by cumulative usage address a documented cognition gap: users score only $2.31$/$5$ on recognising manipulation \parencites[31]{9360054} despite high self-rated confidence, so prompts must be salient. Structural ``natural exit'' points break autoplay's gravitational pull, and staged-progression gamification \parencites[9]{5645864} builds intrinsic motivation better than open-ended streams. Session triggers must be logged as auditable events, mirroring human-in-the-loop grading's dual-logging \parencites[20]{6805789}, as required by AI Law Model event-logging duties \parencites[19.5]{2977039}[228]{2977039} and the anticipated DFA \parencites[390]{7896734}[2248]{6502223}. Recommenders trained on sequential context produce increasingly extreme confidence scores \parencites[4]{8020829} that steer vulnerable users toward harmful communities, making hard session boundaries a technical safety measure.

\subsubsection{Batch Feeds, Digest Modes and Time-Boxed Usage}
\label{sec:batch_feeds_digest_timeboxed} Douyin users average more than 21 sessions daily \parencites[2]{8958037}[3]{8958037}; over half exceed 60 minutes of watch time, concentrated disproportionately among adolescents with weaker self-control. Hyper-personalised feeds correlate \parencites[2]{3303310}[20]{3303310} with $23\%$ higher depressive-symptom incidence. Kostenko's model treats \parencites[119]{2977039}[223]{2977039}[221]{2977039} access to a ``psycho-emotionally safe digital environment'' as an inalienable right, prohibiting algorithms that impose behavioural patterns or form addiction, and mandating that feeds for minors default to ``detox mode.''

A batch or digest mode that compresses sessions without removing manipulative features will not close this gap. Literacy surveys document a ``second-level divide'': users score between $4.40$ and $4.62$ on basic-task confidence \parencites[10]{9360054} yet only $2.31$ out of $5$ on recognising manipulation. Batching must therefore pair with transparent session-boundary signals, consistent with AI Law Model explainability requirements \parencites[68]{2977039}[69]{2977039}. Batch and digest configurations backed by event logs \parencites[242]{2977039} provide the evidential record compliance requires; this report classifies them as ``time-bounded, audit-compatible sessions.''

Audit checks include whether maximum session lengths are documented, whether digest alternatives are surfaced as first-class options, and whether choices are logged for regulator access under DSA Article 40 \parencites[59]{2702055}[19]{3195574}. Immersion research \parencites[8]{9529573} adds a further caution: strong platform immersion can prolong disengagement from offline context, so explicit session boundaries protect better than informative nudges alone. Legibility, default activation, and auditability together determine whether a batch mode satisfies regulatory expectations \parencites[148]{2977039}[390]{7896734}.

\subsubsection{User Control Dashboards for Algorithmic Preferences}
\label{sec:user_control_dashboards} Kostenko's model specifies \parencites[22]{2977039} three obligations: users must be able to know AI's impact on them, obtain explanations, and retain control; every system must disclose its algorithmic nature and support auditing. The Algorithmic Impact Navigator frames \parencites[2]{2742041} dashboards as the operational surface where profiling and recommendation-parameter choices become visible; without a compliant dashboard, those governance obligations remain administrative.

The AI Act requires transparent, auditable and bias-free algorithms with user opt-out of personalisation; the DSA obliges \parencites[8]{5142478}[12]{6502223} very large platforms to offer at least one non-profiled feed. A compliant dashboard must surface the non-profiled mode as a first-class option, logged for audit \parencites[12]{6502223}[21]{3195574}. XAI-based explanation panels stating which factors drove a recommendation, paired with diversity-weighting controls \parencites[3]{3303310}[8]{5142478}, translate bias metrics into user-facing levers; the non-profiled mode must be the default for minors.

A dashboard without an auditable backend is theatre. ISO/IEC AI risk-management standards and the AI Law Model's event-logging duty require \parencites[2]{8589961}[44]{2977039}[135]{2977039} that inputs, calculations and outputs remain traceable and appeal-accessible. HITL grading offers a practical model: instructors see SHAP-based explanations, can override outputs, and every action is dual-logged \parencites[20]{6805789}{2055956}, measurably increasing student trust. Recommender dashboards that separate automated from human-curated content and log overrides achieve the same traceability \parencites[22]{2977039}[61--62]{2977039}.

\subsection{Recommender System Redesign for Wellbeing}

\subsubsection{Diversity-Oriented Ranking and Content Pluralism}
\label{sec:diversity_ranking_pluralism} YouTube's watch-time optimisation amplified sensationalist vaccine content; Facebook disproportionately surfaced far-right political material \parencites[5]{5142478}[29]{5142478}[7]{5142478}. These are documented case findings, established by peer-reviewed UX and neurocognitive research, about how engagement-maximising personalisation operates in practice.

Chitra et al. propose \parencites[16]{6502223} ``value-aligned objective functions'' as a technical response: accuracy weights for fact-checked sources, intra-session diversity requirements, and engagement ceilings that cap amplification of any single item. Gerr argues \parencites[8]{5142478}[28]{5142478} that benchmarks such as Shannon entropy and Simpson's diversity index should be binding design requirements, because DSA reporting compliance alone has not reduced engagement-driven bias on YouTube or TikTok. Deldjoo et al. add precision \parencites[28]{6577872}[30]{6577872} by distinguishing exposure (visibility) from effectiveness (relevance-match) and proposing the Gini index for catalogue concentration as a pluralism indicator suitable for an audit toolkit. Voluntary disclosure of diversity metrics without binding targets has not moved platform behaviour in the reviewed literature.

The trade-offs are real and should not be minimised. Promoting long-tail items without quality controls risks irrelevant recommendations \parencites[28]{6577872}. Capping virality will reduce short-term engagement revenue, placing diversity-oriented ranking in structural conflict with current platform incentive systems \parencites[14]{6502223}. This report treats that as a deliberate rebalancing: engagement metrics may worsen while wellbeing outcomes improve, including reduced depressive-symptom incidence associated with hyper-personalised feeds \parencites[20]{3303310}. Annual independent audits assessing political amplification asymmetries \parencites[16]{6502223}, together with AI Law Model requirements for technical passports and public registers for high-risk systems \parencites[19.5]{2977039}, provide the enforcement architecture. Entropy, Simpson-index and group-utility measures are the primary audit evidence for whether recommenders shift toward pluralistic curation under the anticipated Digital Fairness regime.

\subsubsection{Downranking of Harmful and Sensational Content}
\label{sec:downranking_harmful_sensational} Internal Facebook documents assessed \parencites[2244]{6502223} the platform's engagement algorithm as a ``significant contributor'' to political polarisation, with $64\%$ of extremist group joins traced to recommendation prompts; TikTok accounts signalling interest in weight loss or politics are rapidly steered toward eating-disorder and extremist content. Instagram's ``Explore'' page increased \parencites[3]{3303310}[20]{3303310} exposure to body-negative imagery by $23\%$ relative to chronological feeds, and algorithmically prioritised content shows $31\%$ higher negative affect than baseline. This report's taxonomy treats content combining high negative affect, polarising framing and documented mental-health harm as a prime downranking target. These findings reflect optimisation outcomes in a system designed to maximise engagement, not edge-case failures.

Three technical levers from Chitra et al.'s value-aligned objective functions address this directly. Accuracy weighting boosts fact-checked sources and penalises repeat misinformation patterns. Diversity incentives penalise ideologically narrow sequences. Engagement ceilings cap virality premiums that reward emotional escalation \parencites[16]{6502223}. Table~\ref{tab:downranking_levers} maps these levers to audit focus. \begin{table}[H] \centering \caption{Working mapping: downranking levers for harmful and sensational content} \label{tab:downranking_levers} \begin{tabular}{p{3.3cm}p{4.8cm}p{4.3cm}} \hline Lever & Mechanism & Audit focus \\ \hline Accuracy weighting & Boosts fact checked sources, penalises repeat misinformation & Evidence of source level weights; fact checker integration \parencites[2252]{6502223} \\ Diversity incentives & Penalises ideologically narrow sequences & Diversity metrics (entropy, Simpson index) over recommendation sessions \parencites[2252]{6502223}[6]{5142478} \\ Engagement ceilings & Caps exposure of single items & Detection of hard caps on impressions and watch time per item \parencites[2252]{6502223} \\ Explainable ranking & User facing reasons for each recommendation & Dashboards stating factors such as prior viewing or trending status \parencites[2252]{6502223}[6]{5142478} \\ Mandatory audits & External checks of amplification asymmetries and pluralism & Published audit summaries assessing political amplification and diversity \parencites[2252]{6502223} \\ \hline \end{tabular} \end{table} Chitra et al. accept that downranking inflammatory content will reduce engagement and advertising precision, but treat this trade as necessary given structural misalignment between platform revenue models and democratic values; they call for mandatory annual independent audits with published methodologies and findings \parencites[2250]{6502223}[2252]{6502223}. Platforms failing required auditing risk losing liability protection under conditional safe-harbour regimes \parencites[2253]{6502223}. The AI Law Model prohibits systems that manipulate behaviour or cause digital addiction outright and recognises a right to a psycho-emotionally safe digital environment \parencites[61--62]{2977039}[119]{2977039}.

Scenario-based probes using DSA Article 40 data-access rights \parencites[21]{3195574} can test recommender behaviour under specific risk narratives. Diversity and distributional fairness metrics \parencites[5]{5142478}[27]{5142478} applied to curated datasets offer a second evidential channel. A pro-social pivot must demonstrate value-aligned objectives \parencites[2252]{6502223}[3]{3303310}, measurable affective-balance improvement, and independent audit trails showing sustained downranking of mental-health-harm content. Without a published audit trail, downranking claims are not independently verifiable.

\subsubsection{Explainable Recommendations and User Feedback Loops}
\label{sec:explainable_recos_feedback_loops} Explainability becomes tangible when interfaces show users why items appear and how human oversight operates. HITL grading offers a tested template: feedback is tagged \parencites[15]{6805789}[20]{6805789} ``AI generated,'' ``human reviewed,'' or ``human added,'' with colour coding that measurably increased student trust when instructors visibly modified machine scores. Recommender feeds can apply the same three-way labelling (``algorithmically suggested,'' ``editorially curated,'' ``friend recommended'') with moderator overrides \parencites[11]{6805789}[17]{6805789} dual-logged by timestamp and identifier; without that override trail, the label is disclosure without accountability.

Kostenko's transparency principle requires \parencites[67]{2977039}[68]{2977039} that architecture, logic and purpose be ``understood, accessible, verifiable, and explanatory'' at every lifecycle stage. A ``why this post'' panel linked to an appeal channel implements this directly. The Algorithmic Impact Navigator calls \parencites[2]{2742041} for continuous stakeholder engagement and transparent records, mirrored in HITL grading's guarantee that appeals are reviewed by human faculty with full audit-trail disclosure \parencites[15]{6805789}.

SHAP and LIME make black-box models inspectable and support AI Act/GDPR compliance, but regulators still require \parencites[11]{7461574}[9]{7461574} independent audits and external oversight committees to reach the model itself. An audit toolkit can require XAI pipelines and aggregated explanation summaries for vetted researchers, consistent with DSA data-access provisions \parencites[19]{3195574}[67]{2977039}. Hu's Algorithmic Impact Assessment Scale (AIAS-5) quantifies \parencites[20]{3303310}[19]{3303310} mental-health risk through emotional volatility and compulsive-usage dimensions; combined with explainable interfaces and human-validated overrides, these define a concrete pro-social pathway to dashboards that explain and track impact, feeding regulator channels capable of triggering model throttling when harm indicators rise.

\subsection{Alternative Engagement and Monetization Models}

\subsubsection{Subscription and Freemium Revenue Streams}
\label{sec:subscription_freemium_streams} Subscription and freemium designs become ethically and legally charged when dark commercial patterns enter the picture. An audit of Indian digital platforms found \parencites[2]{6295677} roughly $87$ percent used at least one dark pattern; interface interference appeared in $70$ percent of cases and obstruction in $63$ percent, with $55$ percent of consumers reporting unexpected charges and only $39$ percent recognising the term ``dark patterns,'' signalling a sharp information asymmetry. The ``roach motel'' is its clearest expression: sign-up is frictionless \parencites[4]{6295677} while cancellation is deliberately complex, eroding autonomy and informed consent.

Liu and Keshavarz label this the ``conversion--trust tradeoff.'' Dark commercial patterns raise immediate compliance but corrode the fairness, autonomy and trust that drive long-term loyalty. Their typology covers five categories \parencites[33]{6087241}: information distortion, interface asymmetry, social and temporal pressure, choice obstruction, and forced action; aggressive implementations carry a documented relationship risk \parencites[35]{6087241}, damaging brand relationships and inviting regulatory scrutiny, with no evidence that short-term conversion gains persist once consumer trust declines.

DSA and AI Act regimes are operationalised in part through the European Centre for Algorithmic Transparency (ECAT), whose remit includes promoting fair algorithmic practices across the platforms it monitors \parencites[3]{5142478}. Chitra et al.'s engagement-ceiling and diversity-incentive requirements \parencites[2252]{6502223} are stated for free-tier feeds only; the source literature does not address paid tiers. Applying the same requirements to paid tiers is an extension made here, not a claim already in the source. It follows from the same regulatory logic: relaxing diversity constraints for paying users would undermine the intent of both DSA and the anticipated DFA.

This report classifies subscription and freemium designs relying on information distortion, interface asymmetry or choice obstruction under Code Category C (``Prohibited or Critical Risk Nudges'') \parencites[33]{6087241}[4]{6295677}. Platforms with clear pricing and symmetric sign-up/cancellation flows fall into lower-risk categories, pending DFA and ECAT interpretive guidance \parencites[3]{5142478}. Audit dimensions include classification against the five dark-pattern categories, measurement of cancellation friction via observable UX steps, and assessment of whether premium tiers modify recommendation objectives against diversity and engagement-ceiling standards. That last dimension extends Chitra et al.'s engagement-ceiling framework \parencites[2252]{6502223} to the premium-tier case; the cited literature itself stops short of making that claim.

\subsubsection{Micro-Learning and Positive Habit Loops}
\label{sec:micro_learning_positive_habits} The core design shift is from sessions optimised for time-on-platform to bounded sequences that build competence. Gamified education platforms such as Legends of Learning, Kahoot! and GradeCraft already demonstrate this: when gamification is embedded systemically (clear tasks, defined stages, diagnostic feedback), personalised configurations show \parencites[9]{5645864} $35$--$40$ percent higher engagement and $25$--$30$ percent better learning outcomes than generic designs.

Blooket's token economy illustrates a workable micro-learning arc: compact question batches, instant feedback, and token allocation toward a narrative goal act as a ``currency for creative agency'' rather than shallow point accrual. Rewards must sit in goal-oriented, context-sensitive sequences \parencites[7]{5645864}; punitive feedback or vague contingencies generate technostress that undermines intrinsic motivation. HITL grading offers a comparable template: each assignment moves through AI suggestions (``automated''), instructor comments (``final''), and dual score logs \parencites[5]{6805789}[9]{6805789}[21]{6805789}, with over $80$ percent of students valuing explicit AI-involvement transparency.

Mapped onto recommenders, positive habit loops would define finite content batches around user-relevant goals, with explicit completion signals replacing open-ended scrolling. Kostenko's AI Law Model requires \parencites[25]{2977039} that session boundaries and reward structures be auditable. An aligned audit toolkit requires evidence of finite modules with natural exit points and event logs documenting session start/end and human-oversight intervention \parencites[21]{6805789}[25]{2977039}.

Self-regulation alone is insufficient. Even where firms like Google and Meta maintain internal AI-ethics divisions, transparency is limited, reinforcing calls \parencites[1753]{7461574} for independent audits aligned with the AI Act and OECD AI Principles. Structured tasks and predictable routines should remain user-configurable for batch size, pacing and notification frequency rather than imposed as a fixed template; explainable-AI techniques such as SHAP and LIME \parencites[66]{2855351} offer one plausible route to periodic fairness audits examining differential impacts of engagement loops across demographic and psychological groups, though no cited source yet applies them to this specific use case.

\subsubsection{Advertiser Alignment with Digital Wellbeing}
\label{sec:advertiser_alignment_digital_wellbeing}
Standard advertising metrics track the same dynamics regulators now flag as harmful. Impressions, click-through rates and dwell time are sensitive to the same filter-bubble and high-affect amplification patterns that generate a $23\%$ higher depressive-symptom incidence in hyper-personalised feeds \parencites[20]{3303310}; for advertisers evaluating inventory on those metrics alone, brand-safety exposure is already built in.

The regulatory pressure compounds this. The DFA treats ``addictive product designs and unfair personalisation'' as unfair commercial practices alongside DSA/DMA dark-pattern provisions \parencites[390]{7896734}. Filter-bubble dynamics documented earlier in this report depress adolescent self-esteem \parencites[3]{3303310}, a wellbeing cost advertisers inherit as brand-safety exposure. Kostenko's AI Law Model adds categorical prohibitions \parencites[217]{2977039} on psychotype- or emotion-based targeted advertising to minors. Algorithmic systems can also reproduce harmful stereotypes \parencites[4]{3278906} by skewing outputs toward demeaning representations of marginalised groups.

Value-aligned objective functions offer a specific alternative. Chitra et al. propose weighting fact-checked sources, enforcing intra-session diversity and capping per-item engagement \parencites[2252]{6502223}; applied to advertising inventory, the same objective function trades some short-term reach for association with less polarising, audited inventory. That application is made here, not argued in the source itself. Advertisers could use mandatory third-party audits to evaluate inventory against a harm-risk index such as the AIAS-5 scale \parencites[20]{3303310}, instead of relying on platform self-declarations.

Practically, this means bidding only on inventory subject to published third-party audits with engagement ceilings and diversity incentives, and prioritising placements that disclose predicted psychological impact rather than burying it in terms of service, in the spirit of the ``digital nutrition label'' model \parencites[20]{3303310}. Campaign teams rewarded for reduced association with high-risk content, alongside standard click-through metrics, would be better positioned to act once the DFA enters force; that governance recommendation is made here, with no equivalent finding in the cited audit or dark-pattern literature.

\subsection{Directional Commercial Viability Analysis}

\subsubsection{Cited Figures versus Modelled Assumptions}
\label{sec:commercial_viability_analysis} A pro-social pivot survives a finance committee only if the revenue case is stated honestly. Three figures here are directly cited from published sources. Roughly $87\%$ of major platforms audited in India deploy at least one dark pattern \parencites[2]{6295677}[4]{6087241}, and $47\%$ of users report difficulty cancelling because of obstructive exit flows. Separately, $81\%$ of platforms self-declare ``dark pattern free'' while complaint resolution sits at only $28\%$ \parencites[12]{6295677}. Engagement-optimised feeds correlate \parencites[20]{3303310} with a $23\%$ higher incidence of depressive symptoms, appearing in churn research as a driver of voluntary abandonment. These three figures set the factual baseline; everything that follows is modelled reasoning.

None of these figures translates directly into ARPU, LTV or subscription conversion without assumptions, stated explicitly below.

Where obstructive cancellation currently retains some share of the $47\%$ of users who report wanting to leave, a redesigned symmetric cancellation flow would be expected to increase voluntary churn in the short term. No cited churn-elasticity figure exists for social recommenders, so the honest modelling position is to treat this as a near-term revenue cost weighed against DFA-era fine exposure \parencites[390]{7896734}. Liu and Keshavarz's conversion--trust tradeoff literature \parencites[37]{6087241} indicates dark-pattern removal costs short-term conversion but protects long-term retention and trust-driven referral, without providing a point estimate this report could cite as fact.

On the advertiser side, bidding only on audited, engagement-ceiling-compliant inventory forgoes some volume-driven reach, a real and near-term cost, but reduces brand-safety exposure to the polarising and body-negative content correlated with the cited $23\%$/$31\%$ harm statistics \parencites[3]{3303310}. The compliance-driven redesign trades a bounded, front-loaded conversion cost for a reduction in regulatory and reputational exposure that the $87\%$/$81\%$/$28\%$ figures show is already material before any new legislation takes effect. A fully quantified model would require primary access to a specific platform's cohort data; this report flags that as future empirical work.

\subsubsection{Illustrative ARPU/LTV Scenario Model}
The preceding argument can be sharpened into a scenario model. Every figure below is a modelled assumption indexed to the report's own baseline; the point is to make the reasoning auditable, not to forecast a real platform's financials. Each row states its modelling basis so a reader can substitute their own assumptions and recompute the outcome.

\begin{table}[H]
\centering
\caption{Illustrative ARPU/LTV Scenario Model for the Pro-Social Pivot (Modelled, Not Cited)}
\label{tab:arpu_ltv_scenario}
\begin{tabular}{p{4.3cm}p{3.8cm}p{5.4cm}}
\hline
Scenario Variable & Modelled Range & Basis \\
\hline
Baseline monthly ARPU (indexed to 100) & 100 (index) & Illustrative reference unit; not a claimed figure for any named platform \\
Short-term conversion loss from symmetric cancellation flow & $-5\%$ to $-12\%$ of baseline ARPU & Modelled against the $47\%$ stated cancellation-difficulty figure and the conversion-trust tradeoff in Liu and Keshavarz \parencites[37]{6087241} \\
Medium-term LTV uplift from trust-driven retention (12-18 months) & $+8\%$ to $+15\%$ LTV & Modelled from the same conversion-trust tradeoff literature \parencites[37]{6087241}; the range is wide because no churn-elasticity point estimate exists for social recommenders \\
DFA-era regulatory fine exposure avoided per enforcement cycle & $1\%$ to $6\%$ of global turnover & Modelled on DSA/AI-Act-era fine-severity bands, anchored to the DFA discussion in Sanikidze \parencites[390]{7896734} and the third-party-audit accountability framework in Chitra et al. \parencites[2252]{6502223} \\
Net directional effect (12-24 month horizon) & Positive once fine-avoidance plus retention uplift exceed the short-term conversion loss & Break-even timing depends on the unmeasured churn-elasticity figure; a testable hypothesis, not a forecast \\
\hline
\end{tabular}
\end{table}

The break-even logic holds under pessimistic assumptions: taking the low end of the retention-uplift range and the high end of the conversion-loss range together, a single DFA-era enforcement action large enough to trigger the modelled fine band would still offset several months of foregone conversion. That asymmetry is the useful output here. What this report contributes is the structure of the calculation; a platform-specific version would swap every modelled row for the organisation's own cohort data.

\section{Critical Reflection, Limitations and Conclusion}

\subsection{Critical Reflection}
\label{sec:critical_reflection} This report's central limitation is methodological and was disclosed from the outset. Every claim about TikTok, Instagram, LinkedIn and comparable platforms rests on secondary sources: regulatory filings, published UX and neurocognitive research, and advocacy-group audits. No claim rests on proprietary ranking code \parencites[6]{6502223}[244]{2977039}, training data or internal telemetry. That constraint is defensible given the qualitative, secondary-data scope defined in Section 1. The audit toolkit can classify externally observable design patterns (autoplay defaults, notification frequency, cancellation friction) with reasonable confidence, while claims about exact optimisation objectives or the causal weight of any single feature remain inferential, triangulated across studies rather than verified against a single authoritative source. Triangulation is a legitimate method; it is not a substitute for primary access.

Coverage across platforms is uneven, reflecting the underlying evidence base rather than selective attention. TikTok and Instagram dominate the audited literature; LinkedIn is genuinely under-studied in the dark-pattern and algorithmic-addiction literature consulted here. Section 4.1.3 addresses this by flagging LinkedIn claims as illustrative extrapolation from adjacent platforms rather than fabricating platform-specific statistics, consistent with the Integrity Rule that every claim traces to a citation or is explicitly marked illustrative. Any organisation applying this toolkit to a professional-networking product should treat the LinkedIn-specific sections as a starting hypothesis requiring primary validation. That gap is an honest acknowledgment, not a structural weakness.

The working coding taxonomy (Code A interface mechanics, Code B psychological triggers, Code C legal risk tiers) is this synthesis's own construction, not an official classification under the DSA, AI Act or the still-unpublished Digital Fairness Act. Risk-tier assignments should be read as defensible interpretations of regulatory intent, not binding legal determinations. Because the DFA remains a policy trajectory rather than adopted law at the time of writing, several forward-looking claims about how addictive design will be regulated are necessarily provisional. Where article or clause numbers eventually diverge from current Commission communications, the normative arguments about autonomy, transparency and duty of care should still hold; they are grounded in the DSA and AI Act as already-adopted anchor instruments.

The commercial-viability reasoning in Section 5 raises a distinct concern. It mixes cited figures with modelled assumptions about how those figures translate into subscription conversion or advertiser revenue. These two categories carry different evidentiary weight and should not be read as equally certain. The audit toolkit is best used as a structured starting point for platform-specific investigation, not a substitute for it.

\subsection{Recommended Validation Methodology}
\label{sec:recommended_validation_methodology} One analyst coded the entire Code A/B/C matrix from secondary sources, with a single-evaluator walkthrough for LinkedIn in Section 4.1.3. That is not the same as independent expert validation, and this report does not pretend otherwise. The protocol below is narrow on purpose: it names specific participant groups, a specific instrument, and specific pass/fail thresholds, because "further research is needed" carries no operational weight as a recommendation.

A useful validation round would recruit 8 to 12 participants across three groups, each with direct exposure to one of the matrix's code families. Compliance officers or legal counsel familiar with the DSA and AI Act would stress-test Code C legal-risk-tier classifications. UX or product designers with production experience on recommender-driven interfaces would stress-test Code A interface-mechanics classifications. Researchers or practitioners in behavioural science or AI ethics would stress-test Code B psychological-trigger classifications. Using all three groups checks whether each code family actually works for the practitioners it is meant to serve; a compliance officer struggling with Code B is as informative a result as strong agreement on Code C, and this triangulation logic matches how the audit methodology itself was built \parencites[16]{8200246}[19]{3195574}.

The instrument itself would be a semi-structured interview. Each participant reviews five to eight already-coded audit excerpts, with the assigned Code A/B/C labels visible, and is asked per excerpt whether they would assign the same code, what they would assign instead if not, and how confident they are in that judgement. A short post-interview Likert-scale questionnaire would then capture how applicable the participant found the toolkit to their own compliance or design workflow, giving at least partially comparable results across participants rather than open-ended impressions alone.

Two criteria would separate a successful validation round from an inconclusive one. First: inter-rater agreement between participants and the original coding above a threshold set in advance, not chosen after seeing the results: a Cohen's kappa or percentage-agreement figure. Second: a majority of participants rating the toolkit as practically applicable to their own workflow, not merely conceptually sound. Falling short of either would be a real finding, not a failure to explain away. It would show which code family, or which platform's coding, needs revision before the toolkit goes into operational use. Naming a protocol this specific is the point: it can be picked up and run by a future researcher without further design work, rather than sitting as an unaddressed gap.

\subsection{Concluding Strategic Outlook}
\label{sec:concluding_strategic_outlook} Engagement-optimised social recommenders combine interface affordances, profiling practices and business incentives in ways that materially affect user autonomy, attention allocation and distributional outcomes. Endless feeds, autoplay, high-frequency re-engagement prompts, asymmetric defaults and opaque opt-out channels create feedback loops the literature \parencites[2]{3303310}[2]{8958037} associates with compulsive checking, degraded wellbeing and narrowed exposure, with neurocognitive, behavioural and UX evidence consistently linking these design choices to reduced self-regulation and elevated affective volatility for vulnerable groups. European legal and model-law frameworks now treat such configurations as regulatory risk rather than acceptable product differentiation. Obligations for transparency, non-profiling feed options, tamper-evident logs and age-sensitive defaults are specified \parencites[12]{6502223}[61--62]{2977039}; where those artefacts are absent or inadequate, governance failure becomes a signal of unacceptable systemic risk under the DSA and AI Act. In practice, a platform cannot redesign its ranking logic for regulatory compliance without also changing what its product team ships.

Compliance readiness and commercial resilience now point in the same direction. The pro-social pivot strategy in Section 5 shows that finite-session design, user control dashboards, diversity-aware ranking and subscription/freemium models decoupled from dark-pattern retention tactics are plausible routes to sustaining engagement and revenue once the current attention-maximising model faces DFA-level scrutiny \parencites[390]{7896734}. Operationally, a practicable audit and remediation pathway should focus on observable, auditable artefacts: documented algorithmic passports, event logs with human-override trails, user control dashboards exposing non-profiled modes, finite session mechanisms and batch or digest alternatives, explainable recommendation panels, and diversity-aware ranking objectives with downranking levers for sensational content (see Section 5). Embedding these controls \parencites[16]{6502223}[244]{2977039} into CI/CD gates, independent third-party audits and governance cycles creates traceable evidence for regulators and civil society while enabling product teams to iterate on safer objective functions without waiting for finalised DFA text. Without those artefacts, regulatory intent has no way to reach the codebase.

Priority empirical work beyond this report should centre on scenario-based probes, representative trace releases for vetted researchers, and longitudinal evaluation of wellbeing outcomes after deploying time-boxed, diversity-oriented interventions. LinkedIn is the clearest gap: the evidence base there remains thin, and primary research would close a documented hole. Each of these steps feeds directly into audit practice: scenario probes generate the DSA Article 40 datasets that external researchers need; trace releases let regulatory bodies verify downranking claims independently; and longitudinal data would finally answer whether the $23\%$ depressive-symptom reduction correlated with chronological feeds is sustained when platforms move from unlimited engagement optimisation to the bounded, wellbeing-oriented designs described in Section 5 \parencites[119]{2977039}[390]{7896734}.

\newpage
\printbibliography

\newpage
\section*{List of Appendices}
\addcontentsline{toc}{section}{List of Appendices}
\begin{itemize}
\item Appendix A: Capstone Project Exposé
\item Appendix B: AI Use Declaration
\end{itemize}

\newpage
\section{Appendix A: Capstone Project Expos\'e}
\label{sec:appendix_a_expose}

\textbf{Name:} Mohit Nirwan

\textbf{Matriculation Number:} IU14087991

\textbf{Programme:} Master of Business Administration (IU MBA 90 ECTS / LSBU Dual-Award)

\subsection*{1. Preliminary Title}
Algorithmic Addiction in Social Recommender Systems - Development of an Ethical AI Audit and Pro-Social Pivot Strategy under the EU Digital Fairness Act

\subsection*{2. Introduction: Problem Description and Objective}
Social media platforms like TikTok, Instagram, and LinkedIn optimize their recommendation engines primarily for screen time. To hit those numbers, the design leans on infinite scroll, variable reward schedules, and content loops built to provoke outrage. All three exploit the same cognitive weak spots. For example, TikTok's auto-playing feed and pull-to-refresh mechanics act as textbook variable reward schedules. This ad-driven model is highly monetizable, but it's wearing users down. Younger users in particular are showing signs of digital burnout and cognitive overload. That reliance on engagement-hacking has now hit a regulatory wall. The European Commission's upcoming Digital Fairness Act (targeted for Q4 2026) explicitly aims to ban deceptive design and addictive algorithmic patterns. If platforms do not comply, they face revenue-based fines.

Right now, most executive teams have no real way to audit their own algorithms, let alone fix them without tanking retention. My Capstone Project turns this into a consulting report with a specific output: a roadmap for moving away from ad-supported dark patterns toward what I'm calling Pro-Social AI. That means testing whether subscription models and empathetic interaction design can replace outrage-driven engagement without sinking revenue. The goal is to answer a specific Research Question: How can enterprise social platforms redesign their AI recommender systems to comply with upcoming EU Digital Fairness standards while protecting core user engagement and commercial viability?

\subsection*{3. Brief Theoretical Foundation}
The analysis is anchored in Behavioural Economics, specifically Thaler and Sunstein's (2008) choice architecture and the taxonomy of dark patterns mapped by Mathur et al. (2019). Two other strands feed into this. Self-Determination Theory (Deci \& Ryan, 2000) explains why users are vulnerable in the first place. More recent work on social-platform dark patterns (Mildner et al., 2023; Wiesb\"ock et al., 2026) shows how that vulnerability gets designed into the product. A clear limit to acknowledge is that behavioural economics often assumes users can be "nudged" rationally, which breaks down in highly addictive, frictionless environments. Still, it provides the best diagnostic baseline.

\subsection*{4. Research Methodology (Method and Analysis)}
The recommendations rest on application-specific secondary data. I will gather three kinds of secondary data: the regulatory texts themselves (the DSA and the 2026 DFA drafts), public UI/UX audits of major platforms, and consumer advocacy submissions on digital burnout going back to 2021. I will code these inputs using qualitative thematic analysis to map specific algorithmic design patterns against legal risk categories. One limitation I will have to manage: secondary data regarding proprietary algorithms relies on reverse-engineered audits rather than open internal codebases. I can bridge this gap by focusing on the observable user interface outputs rather than the underlying black-box code.

\subsection*{5. Planned Structure of the Consulting Report}
The following outline aligns directly with the IU consulting report requirements:
\begin{itemize}
\item 1. Cover Page
\item 2. Summary
\item 3. Problem description and environment: Algorithmic engagement optimization, addictive design, digital burnout, and the legislative threat of the EU Digital Fairness Act.
\item 4. Consulting function
\begin{itemize}
\item 4.1. Objective and challenges: Harmonizing legal compliance and pro-social features with platform business models.
\item 4.2. Method: Qualitative thematic coding of regulatory standards, consumer burnout data, and UI design patterns.
\item 4.3. Analysis: Comparative audit of current platform recommendation loops against proposed legal boundaries.
\item 4.4. Recommendations: An audit tool for compliance and a strategic pivot roadmap toward Pro-Social AI (e.g., subscription models, micro-learning loops).
\end{itemize}
\item 5. Critical reflection: Concluding discussion of process limitations.
\item 6. Literature
\item 7. Appendix: (Including this expos\'e for the LSBU award assessment)
\end{itemize}

\subsection*{6. Initial Bibliography}
\begin{itemize}
\item Deci, E. L., \& Ryan, R. M. (2000). The "what" and "why" of goal pursuits: Human needs and the self-determination of behaviour. \textit{Psychological Inquiry, 11}(4), 227-268.
\item European Commission. (2022). Regulation (EU) 2022/2065 (Digital Services Act). \textit{Official Journal of the European Union}.
\item European Parliament. (2026). Digital Fairness Act: Protecting our democracy, upholding our values. \textit{Legislative Train Schedule} (Q4 2026 Initiative).
\item Mathur, A., et al. (2019). Dark Patterns at Scale: Findings from a Crawl of 11K Shopping Websites. \textit{Proceedings of the ACM on Human-Computer Interaction, 3}(CSCW), 1-32.
\item Mildner, T., Savino, G. L., Doyle, P. R., Cowan, B. R., \& Malaka, R. (2023). About Engaging and Governing Strategies: A Thematic Analysis of Dark Patterns in Social Networking Services. \textit{Proceedings of the 2023 CHI Conference on Human Factors in Computing Systems}, 1-17.
\item Thaler, R. H., \& Sunstein, C. R. (2008). \textit{Nudge: Improving decisions about health, wealth, and happiness}. Yale University Press.
\item Wiesb\"ock, L., Wang, N., Reitzer, P., B\u{a}noiu, I., Beltzung, L., \& Lehner, P. (2026). ADDICT: Addictive Design in Digital Consumer Technologies. Verlag Arbeiterkammer Wien.
\end{itemize}

\newpage
\section{Appendix B: AI Use Declaration}
\label{sec:appendix_b_ai_declaration}
This thesis was prepared with the assistance of large language model tools that were used to generate, edit, and refine portions of the text under the direct supervision of the author. Such tools were employed to help draft language, suggest rephrasings, and identify stylistic or organizational improvements; all outputs from these tools were reviewed, modified as necessary, and incorporated at the author's discretion.

The author independently verified and integrated all source material, checked factual statements, and ensured that citations, data, and interpretations are accurate and appropriate for an MBA capstone thesis. Where text proposed by a language model was used, the author confirmed its conformity with academic standards, adapted the content to the thesis argument, and ensured proper attribution of ideas and evidence.

Ultimately, responsibility for the content, analysis, argumentation, and academic integrity of this work rests solely with the author. Any errors, omissions, or interpretations in the final document are the author's own.

\end{document}
